% TalTech Doctoral Thesis
% Agent-Enhanced Platform Engineering
% Siim Vene, 2026

\documentclass{TTUPhD}

% Thesis metadata
\thesistitle{Agent-Enhanced Platform Engineering: Trust Frameworks and Patterns for Observability and Incident Response}
\thesistitleet{Agentidega täiustatud platvormiarendus: usaldusraamistikud ja mustrid jälgitavuse ning intsidentidele reageerimise jaoks}
\thesisauthor{Siim Vene}
\supervisor{[Supervisor Name], PhD, TalTech}
% \cosupervisor{[Co-supervisor Name], PhD, [Institution]}
% \opponent{[Opponent Name], PhD, [Institution]}
% \defensedate{[Date]}
% \issn{XXXX-XXXX}
% \isbn{XXX-XXX-XXXX-XX-X}

\begin{document}

% Front matter
\frontmatter

% Title pages
\maketitlepage
\maketitlepageet
\makereversepage

% Table of contents
\tableofcontents
\listoffigures
\listoftables

% Publications (if article-based thesis)
\begin{publications}
This thesis is a monograph. No prior publications are included.
\end{publications}

% Author contributions
\begin{contributions}
All work presented in this thesis is the original contribution of the author, conducted under the supervision of [Supervisor Name].
\end{contributions}

% Abbreviations
\begin{abbreviations}
\begin{tabular}{ll}
AI & Artificial Intelligence \\
API & Application Programming Interface \\
CI/CD & Continuous Integration/Continuous Deployment \\
DORA & DevOps Research and Assessment \\
IDP & Internal Developer Platform \\
K8s & Kubernetes \\
LLM & Large Language Model \\
MCP & Model Context Protocol \\
MTTR & Mean Time to Recovery \\
NIS2 & Network and Information Security Directive 2 \\
PAM & Privileged Access Management \\
RBAC & Role-Based Access Control \\
RCA & Root Cause Analysis \\
SLO & Service Level Objective \\
SRE & Site Reliability Engineering \\
\end{tabular}
\end{abbreviations}

% Main matter
\mainmatter

% Part I: Foundations
\part{Foundations}
\chapter{Introduction}
\label{ch:introduction}

% Working draft - S. Vene, 2026-02-08

\section{The Operational Challenge at Scale}
\label{sec:operational-challenge}

Large software organizations operate complex, distributed systems that are increasingly difficult to observe and operate reliably. Platform engineering and Site Reliability Engineering (SRE) practices have emerged to manage this complexity, supported by observability tooling and automation. Despite these advances, operational teams face persistent challenges:

\begin{itemize}
    \item \textbf{Signal overload:} Modern systems generate telemetry volumes that exceed human processing capacity
    \item \textbf{Cognitive load:} Incident diagnosis requires correlating signals across dozens of services and data sources
    \item \textbf{Expertise bottleneck:} Senior engineers are required for complex diagnosis, creating 24/7 coverage challenges
    \item \textbf{Slow resolution:} Root cause analysis consumes 40--60\% of incident time
    \item \textbf{Learning loss:} Incident knowledge doesn't transfer effectively to prevention
\end{itemize}

The DORA 2024 report found that reducing Mean Time to Recovery (MTTR) remains critical for organizational performance, yet operational toil continues to consume platform engineering resources.

\section{The Agent Opportunity}
\label{sec:agent-opportunity}

Recent advances in large language models (LLMs) and agentic AI systems create opportunities to move beyond statistical AIOps approaches. Unlike rule-based automation or statistical anomaly detection, AI agents can:

\begin{itemize}
    \item \textbf{Synthesize heterogeneous telemetry:} Correlate logs, metrics, traces, and change events
    \item \textbf{Reason about system behavior:} Form diagnostic hypotheses based on evidence
    \item \textbf{Retrieve institutional knowledge:} Access past incidents, runbooks, and documentation
    \item \textbf{Support or execute remediation:} From suggesting actions to executing approved runbooks
\end{itemize}

Raza et al. \citep{raza2025} define Agentic AI as ``Multi-Agent Systems powered by LLMs that exhibit autonomous planning, tool use, memory retention, and emergent reasoning capabilities, with or without human supervision.'' This capability set aligns precisely with operational challenges: agents excel at pattern recognition, correlation, and knowledge retrieval---the cognitive tasks that consume human responder time.

\section{The Trust Problem}
\label{sec:trust-problem-intro}

However, deploying autonomous agents in production systems creates trust challenges that existing frameworks don't address:

\begin{itemize}
    \item \textbf{Probabilistic behavior:} Unlike deterministic automation, agents may produce different outputs for the same input
    \item \textbf{High-stakes context:} Production incidents have business impact; wrong actions make things worse
    \item \textbf{Limited oversight:} Agents may operate at 3am when human oversight is minimal
    \item \textbf{Unclear boundaries:} When should an agent act autonomously vs. defer to humans?
\end{itemize}

The DORA 2024 report found that 39\% of respondents have little or no trust in AI-generated code. For AI-generated remediation actions in production, trust requirements are even higher.

Recent empirical work confirms these concerns are well-founded. Riddell et al. \citep{riddell2026} tested six LLMs across 48,000 simulated cloud root cause analysis scenarios and identified a taxonomy of 16 distinct reasoning failure modes. LLMs exhibited systematic vulnerabilities: ``stalled'' (reasoning loops), ``biased'' (anchoring on irrelevant signals), and ``confused'' (multi-hop reasoning failures). Critically, these failures are \textit{unpredictable}---the same model that succeeds on one scenario may fail on a structurally similar one. This empirical evidence underscores that trust in agentic systems cannot be assumed; it must be engineered through explicit safeguards.

This creates a fundamental tension: agents are most valuable where they can act autonomously (reducing human toil), but autonomous action requires trust that has not been established. Organizations need frameworks for calibrating appropriate autonomy levels and building trust incrementally.

\section{Research Questions}
\label{sec:research-questions}

This thesis investigates how AI agents can enhance operational aspects of platform engineering while maintaining appropriate trust and governance. The research questions are:

\begin{description}
    \item[RQ1:] How can agentic systems enhance observability by transforming raw telemetry into actionable operational understanding?
    \item[RQ2:] To what extent can agentic reasoning improve incident triage and root cause analysis compared to existing SRE practices?
    \item[RQ3:] Under what conditions can autonomous or semi-autonomous agents safely and effectively perform remediation actions in production systems?
    \item[RQ4:] What trust and governance mechanisms enable organizations to calibrate appropriate autonomy levels for operational agents?
\end{description}

\section{Methodology}
\label{sec:methodology}

This thesis employs Design Science Research \citep{hevner2004,peffers2007}, appropriate for developing artifacts (frameworks, patterns) that solve identified problems.

\subsection{Design Science Approach}
\label{subsec:design-science}

The research cycle:
\begin{enumerate}
    \item \textbf{Problem identification:} Operational challenges in observability, triage, remediation
    \item \textbf{Objective definition:} Trust framework, operational patterns, governance guidance
    \item \textbf{Design and development:} Iterative framework and pattern development
    \item \textbf{Demonstration:} Application to operational scenarios
    \item \textbf{Evaluation:} Practitioner review, constraint stress-testing
    \item \textbf{Communication:} This thesis
\end{enumerate}

\subsection{Grounding in Practice}
\label{subsec:grounding-practice}

The research is grounded in 15 years of platform and infrastructure architecture experience:
\begin{itemize}
    \item Chief Infrastructure Architect at SMIT (government, 100+ teams)
    \item Domain Architect at Swedbank (financial services, regulated environment)
    \item Program Manager at TalTech (teaching, student feedback)
    \item Founder at Kleidia (NIS2 compliance context)
\end{itemize}

This experience shapes problem identification and provides the constraint space against which solutions are tested.

\subsection{Limitations Acknowledged}
\label{subsec:limitations-acknowledged}

This is practitioner-grounded research, not controlled experiments:
\begin{itemize}
    \item Patterns are proposed based on experience and stress-testing, not empirical measurement
    \item The trust framework is grounded in theory \citep{mayer1995} but not empirically validated
    \item Contributions are offered as hypotheses for practice, not proven solutions
\end{itemize}

\section{Contributions}
\label{sec:contributions}

This thesis makes four contributions:

\begin{enumerate}
    \item \textbf{Trust Framework (Chapter~\ref{ch:trust-framework}):} A constraint-based model for calibrating agent autonomy, grounded in organizational trust theory. Given a target autonomy level, the framework specifies minimum safeguards required.

    \item \textbf{Operational Patterns (Chapters~\ref{ch:observability}--\ref{ch:remediation}):} Documented patterns for agent-enhanced observability, incident triage, and remediation, with autonomy levels and implementation guidance.

    \item \textbf{Governance Guidance (Chapter~\ref{ch:governance}):} Enterprise deployment patterns addressing data sovereignty, cost governance, and LLM infrastructure trust.

    \item \textbf{Practitioner Validation (Chapter~\ref{ch:evaluation}):} Honest assessment of limitations and future work required for empirical validation.
\end{enumerate}

\section{Scope and Delimitations}
\label{sec:scope}

\textbf{Included:}
\begin{itemize}
    \item Observability: transforming telemetry into understanding
    \item Incident triage: diagnosis, root cause analysis, hypothesis generation
    \item Remediation: runbook execution, rollback, scaling (with guardrails)
    \item Trust and governance for operational agents
\end{itemize}

\textbf{Excluded:}
\begin{itemize}
    \item Design-time agent applications (architecture review, documentation generation)
    \item Build-time agent applications (code review, test generation)
    \item Foundation model development
    \item Fully autonomous operation without human oversight
    \item Purely statistical anomaly detection (not agentic)
\end{itemize}

\section{Thesis Structure}
\label{sec:structure}

\textbf{Part I: Foundations}
\begin{itemize}
    \item Chapter 1: Introduction (this chapter)
    \item Chapter 2: Background and Literature Review
\end{itemize}

\textbf{Part II: Trust and Governance}
\begin{itemize}
    \item Chapter 3: The Trust Problem
    \item Chapter 4: A Trust Framework for Operational Agents
    \item Chapter 5: Governance at Enterprise Scale
\end{itemize}

\textbf{Part III: Operational Patterns}
\begin{itemize}
    \item Chapter 6: Agent-Enhanced Observability
    \item Chapter 7: Agent-Enhanced Incident Triage
    \item Chapter 8: Agent-Assisted Remediation
\end{itemize}

\textbf{Part IV: Implementation and Evaluation}
\begin{itemize}
    \item Chapter 9: Implementation Approaches
    \item Chapter 10: Evaluation and Limitations
    \item Chapter 11: Conclusion
\end{itemize}

\chapter{Background and Literature Review}
\label{ch:literature-review}

% Working draft - S. Vene, 2026-02-08

\section{Positioning This Review}
\label{sec:positioning-review}

Chapter~\ref{ch:introduction} traced the evolution of platform engineering through four phases: manual operations, Infrastructure as Code, GitOps, and now agent-enhanced systems. This literature review examines the academic and industry research that informs each of these phases, with particular focus on the emerging body of work on AI agents in enterprise contexts.

\section{AI Agent Architectures}
\label{sec:agent-architectures}

\subsection{Defining Agentic AI}
\label{subsec:defining-agentic-ai}

Following Raza et al. \citep{raza2025}, we define Agentic AI as:

\begin{quote}
``Multi-Agent Systems (MAS) powered by LLMs that exhibit autonomous planning, tool use, memory retention, and emergent reasoning capabilities, with or without human supervision.''
\end{quote}

Key characteristics distinguishing agentic AI from traditional automation:
\begin{itemize}
    \item \textbf{Autonomous planning}---decomposing goals into sub-tasks
    \item \textbf{Tool use}---invoking external APIs and systems
    \item \textbf{Memory retention}---maintaining context across interactions
    \item \textbf{Emergent reasoning}---behavior not explicitly programmed
\end{itemize}

\subsection{Multi-Agent System Architectures}
\label{subsec:mas-architectures}

Three dominant paradigms have emerged:

\begin{table}[htbp]
\centering
\caption{Multi-Agent System Paradigms}
\label{tab:mas-paradigms}
\begin{tabular}{lll}
\toprule
\textbf{Paradigm} & \textbf{Representative Framework} & \textbf{Key Characteristics} \\
\midrule
Graph-based & LangGraph & Explicit state machines, deterministic routing, high auditability \\
Role-based & CrewAI & Agents as team members with roles/goals, intuitive mental model \\
Conversational & AutoGen & Agents communicate via messages, emergent coordination \\
\bottomrule
\end{tabular}
\end{table}

Each paradigm makes different tradeoffs between control and flexibility. Graph-based approaches offer maximum control, while conversational approaches offer maximum flexibility.

\subsection{Cognitive Architectures for Language Agents (CoALA)}
\label{subsec:coala}

Sumers et al. \citep{sumers2024} provide a unifying conceptual framework for language-based agents through the Cognitive Architectures for Language Agents (CoALA) model. CoALA structures agent systems via three core components:

\textbf{1. Modular Memory}
\begin{itemize}
    \item Working memory (current context)
    \item Episodic memory (past experiences)
    \item Semantic memory (learned knowledge)
    \item Procedural memory (skills and tools)
\end{itemize}

\textbf{2. Structured Action Spaces}
\begin{itemize}
    \item Internal actions (reasoning, planning, retrieval)
    \item External actions (tool use, environment interaction)
    \item Action selection policies
\end{itemize}

\textbf{3. Decision Processes}
\begin{itemize}
    \item Grounded decision-making
    \item Explicit reasoning traces
    \item Evidence-based action selection
\end{itemize}

\textbf{Relevance to Trust:} CoALA's modular architecture enables fine-grained observability---each memory type and action can be audited independently. This supports the Trust Equation framework (Chapter~\ref{ch:trust-framework}) by providing natural observation points for monitoring agent behavior.

\subsection{Session-Based vs In-Memory Orchestration}
\label{subsec:session-vs-inmemory}

A fourth paradigm, not well-covered in current literature, is \textbf{session-based orchestration} where agents run as isolated sessions with persistent state:

\begin{table}[htbp]
\centering
\caption{Session-Based vs In-Memory Orchestration}
\label{tab:session-vs-inmemory}
\begin{tabular}{lll}
\toprule
\textbf{Property} & \textbf{In-Memory} & \textbf{Session-Based} \\
\midrule
State persistence & Lost on failure & Survives restarts \\
Isolation & Shared process & Separate sessions \\
Auditability & Requires instrumentation & Native (session logs) \\
Latency & Low (ms) & Higher (seconds) \\
Recovery & Complex & Simple (resume session) \\
\bottomrule
\end{tabular}
\end{table}

\textbf{Research gap:} Current surveys focus on in-memory architectures. The tradeoffs of session-based orchestration for trust and safety are underexplored.

\subsection{Standards for Agent Interoperability}
\label{subsec:agent-standards}

\subsubsection{IEEE P3394: Universal Message Format for LLM Agents}

The IEEE P3394 standard (in development, 2024--2026) addresses a critical gap in multi-agent systems: interoperability. As organizations deploy agents from multiple vendors and frameworks, the lack of standardized communication protocols creates integration challenges.

P3394 defines:
\begin{itemize}
    \item \textbf{Universal Message Format (UMF)}---structured JSON/Protocol Buffer schemas for agent-to-agent communication
    \item \textbf{Session management}---standardized handshakes, context passing, and termination
    \item \textbf{Capability advertisement}---agents declare available tools and competencies
    \item \textbf{Error handling}---consistent error codes and recovery protocols
\end{itemize}

\subsubsection{IEEE P3428: Modular Agent Architecture}

Complementing P3394, IEEE P3428 specifies a modular agent architecture supporting:
\begin{itemize}
    \item \textbf{Plug-and-play integration}---agents can be composed without custom adapters
    \item \textbf{Lifecycle management}---standardized states (initializing, ready, executing, suspended, terminated)
    \item \textbf{Adaptive learning interfaces}---how agents update their behavior based on feedback
\end{itemize}

\textbf{Relevance to this thesis:} These emerging standards provide a foundation for the enterprise deployment patterns in Chapter~\ref{ch:implementation}. However, neither standard addresses trust calibration or governance---gaps this thesis aims to fill.

\section{Trust and Trustworthiness in Autonomous Systems}
\label{sec:trust-theory}

\subsection{Defining Trust vs Trustworthiness}
\label{subsec:trust-vs-trustworthiness}

Raza et al. \citep{raza2025} make a critical distinction:

\begin{quote}
``\textbf{Trust} refers to a user's willingness to rely on an AI system, while \textbf{trustworthiness} denotes whether the system consistently behaves in a safe, fair, and predictable manner, thereby deserving that trust.''
\end{quote}

This distinction is crucial for enterprise adoption:
\begin{itemize}
    \item Users may \textbf{trust} a system that is not \textbf{trustworthy} (dangerous)
    \item A \textbf{trustworthy} system may not be \textbf{trusted} (adoption failure)
\end{itemize}

Building trustworthy systems is necessary but not sufficient---organizations must also build appropriate levels of trust.

\subsection{Foundational Trust Theory}
\label{subsec:foundational-trust}

\subsubsection{Organisational Trust: Mayer, Davis \& Schoorman (1995)}

The seminal organisational trust model identifies three antecedents of trust:

\begin{table}[htbp]
\centering
\caption{Mayer et al. Trust Antecedents Applied to Agents}
\label{tab:mayer-trust}
\begin{tabular}{lll}
\toprule
\textbf{Antecedent} & \textbf{Definition} & \textbf{Agent Equivalent} \\
\midrule
Ability & Skills and competencies enabling influence & Task completion accuracy, tool proficiency \\
Benevolence & Belief trustee wants good for trustor & Alignment with operator goals, no hidden agendas \\
Integrity & Adherence to principles trustor finds acceptable & Policy compliance, predictable behavior \\
\bottomrule
\end{tabular}
\end{table}

\begin{quote}
``Trust is the willingness of a party to be vulnerable to the actions of another party based on the expectation that the other will perform a particular action important to the trustor, irrespective of the ability to monitor or control that other party.'' \citep{mayer1995}
\end{quote}

\textbf{Application:} This model explains why \textit{capability alone} is insufficient for agent adoption. Organisations must also perceive that agents are \textit{aligned} (benevolence) and \textit{predictable} (integrity).

\subsubsection{Trust in Automation: Lee \& See (2004)}

Lee \& See's framework is foundational for understanding human-automation trust \citep{lee2004}:

\textbf{Key Concepts:}

\begin{enumerate}
    \item \textbf{Appropriate Reliance}---Trust should match system capability
    \begin{itemize}
        \item \textit{Misuse} (overtrust)---relying when system is inadequate
        \item \textit{Disuse} (undertrust)---failing to rely when system is adequate
        \item \textit{Abuse}---using automation improperly
    \end{itemize}

    \item \textbf{Trust Calibration}---Alignment between trust and trustworthiness
    \begin{itemize}
        \item \textit{Resolution}---sensitivity to changes in trustworthiness
        \item \textit{Specificity}---granularity of trust judgments
    \end{itemize}

    \item \textbf{Bases of Trust:}
    \begin{itemize}
        \item \textit{Performance}---current behavior
        \item \textit{Process}---understanding of mechanisms
        \item \textit{Purpose}---perceived intent/design goals
    \end{itemize}
\end{enumerate}

\textbf{Critical Insight:}
\begin{quote}
``Trust in automation should be designed for appropriate reliance, because misuse and disuse reduce performance and can increase risk.''
\end{quote}

This framing directly motivates the staged autonomy approach: rather than maximising trust, the goal is \textit{calibrating} trust to actual capability.

\subsubsection{Levels of Automation: Parasuraman, Sheridan \& Wickens (2000)}

Parasuraman et al. \citep{parasuraman2000} provide a structured model for automation across four stages:

\begin{table}[htbp]
\centering
\caption{Levels of Automation by Stage}
\label{tab:loa-stages}
\begin{tabular}{lllll}
\toprule
\textbf{Stage} & \textbf{Function} & \textbf{Low Automation} & \textbf{High Automation} \\
\midrule
1 & Information Acquisition & Human gathers & System gathers \\
2 & Information Analysis & Human interprets & System interprets \\
3 & Decision Selection & Human decides & System recommends/decides \\
4 & Action Implementation & Human acts & System acts \\
\bottomrule
\end{tabular}
\end{table}

\textbf{Key Insight:} Automation does not simply replace human work---it \textit{transforms} it. High automation at Stage 4 (action) may increase cognitive load at Stage 3 (decision oversight).

\subsection{The TRiSM Framework}
\label{subsec:trism}

Gartner's AI TRiSM (Trust, Risk, and Security Management) framework, adapted for agentic AI by Raza et al. \citep{raza2025}, comprises four pillars:

\textbf{Pillar 1: Explainability}
\begin{itemize}
    \item Decision provenance---why did the agent take this action?
    \item Chain-of-thought visibility
    \item Cross-agent communication transparency
\end{itemize}

\textbf{Pillar 2: ModelOps}
\begin{itemize}
    \item Agent lifecycle management
    \item Version control for agent configurations
    \item Deployment pipelines for agent updates
    \item Rollback capabilities
\end{itemize}

\textbf{Pillar 3: Security \& Privacy}
\begin{itemize}
    \item Authentication between agents
    \item Data access controls
    \item Prompt injection defense
    \item Privacy-preserving computation
\end{itemize}

\textbf{Pillar 4: Governance}
\begin{itemize}
    \item Policy enforcement
    \item Compliance monitoring
    \item Human oversight mechanisms
    \item Accountability structures
\end{itemize}

\subsection{Risk Taxonomy for Agentic AI}
\label{subsec:risk-taxonomy}

Raza et al. \citep{raza2025} identify unique threat vectors in multi-agent systems:

\begin{table}[htbp]
\centering
\caption{Risk Taxonomy for Agentic AI}
\label{tab:risk-taxonomy}
\begin{tabular}{lll}
\toprule
\textbf{Risk Category} & \textbf{Description} & \textbf{Example} \\
\midrule
Coordination failures & Agents fail to synchronize & Builder and reviewer on different versions \\
Prompt injection & Malicious inputs manipulate behavior & Data contains override instructions \\
Memory poisoning & Corrupted context propagates & Wrong info in shared memory \\
Agent collusion & Agents coordinate against goals & Agents ``agree'' to skip validation \\
Emergent misbehavior & Unexpected behavior from interactions & Reward hacking in multi-agent systems \\
\bottomrule
\end{tabular}
\end{table}

\subsection{Zero Trust Architecture for Agentic Systems}
\label{subsec:zero-trust}

\subsubsection{From Network Security to Agent Security}

Zero Trust Architecture (ZTA), codified in NIST SP 800-207 \citep{nist800207}, has become the dominant security paradigm for enterprise networks. The core principle---``never trust, always verify''---applies directly to agentic AI:

\begin{table}[htbp]
\centering
\caption{Zero Trust Principles Applied to Agents}
\label{tab:zero-trust-agents}
\begin{tabular}{lll}
\toprule
\textbf{ZTA Principle} & \textbf{Network Application} & \textbf{Agent Application} \\
\midrule
Verify explicitly & Authenticate every request & Verify every agent action against policy \\
Least privilege & Minimal network access & Minimal tool/data access per task \\
Assume breach & Segment networks & Isolate agent sessions, limit blast radius \\
\bottomrule
\end{tabular}
\end{table}

\subsubsection{Cloud Security Alliance Agentic Trust Framework}

The CSA's Agentic Trust Framework (2026) extends Zero Trust to autonomous agents, proposing five governance questions:

\begin{enumerate}
    \item \textbf{Identity}---How do we authenticate agents and verify their provenance?
    \item \textbf{Behavior}---How do we monitor and constrain agent actions in real-time?
    \item \textbf{Data}---How do we control what data agents can access and generate?
    \item \textbf{Segmentation}---How do we isolate agents to limit failure propagation?
    \item \textbf{Incident Response}---How do we detect, respond to, and recover from agent misbehavior?
\end{enumerate}

The framework proposes an ``Intern to Principal'' maturity model where agents earn expanded privileges through demonstrated reliability---a concept this thesis develops further in the Trust Equation framework (Chapter~\ref{ch:trust-framework}).

\subsection{Human-AI Trust Calibration}
\label{subsec:trust-calibration}

\subsubsection{The Calibration Problem}

Trust calibration refers to the alignment between a user's trust in a system and the system's actual trustworthiness. Miscalibration manifests in two failure modes:

\begin{itemize}
    \item \textbf{Overtrust}---Users rely on agents for tasks beyond their competence, leading to errors
    \item \textbf{Undertrust}---Users fail to leverage capable agents, negating productivity benefits
\end{itemize}

\subsubsection{Dispositional, Situational, and Learned Trust}

Building on Hoff \& Bashir's (2015) three-layer trust model \citep{hoff2015}:

\begin{table}[htbp]
\centering
\caption{Three-Layer Trust Model}
\label{tab:three-layer-trust}
\begin{tabular}{lll}
\toprule
\textbf{Trust Layer} & \textbf{Description} & \textbf{Calibration Mechanism} \\
\midrule
Dispositional & General tendency to trust automation & Education, organizational culture \\
Situational & Context-specific trust adjustments & Real-time transparency, status displays \\
Learned & Trust developed through experience & Consistent performance, graceful failures \\
\bottomrule
\end{tabular}
\end{table}

\subsubsection{Transparency as Trust Mechanism}

Agent transparency---making reasoning visible to users---has emerged as a key calibration mechanism. Studies show:
\begin{itemize}
    \item Explanations increase trust when agents are competent
    \item Explanations \textit{decrease} trust when agents make errors (appropriate calibration)
    \item The timing and detail level of explanations affect calibration quality
\end{itemize}

\textbf{Research gap:} Most transparency research examines single-agent systems. How transparency works in multi-agent environments---where users must trust coordination, not just individual decisions---remains underexplored.

\subsection{Measuring Trust in Automation}
\label{subsec:measuring-trust}

\subsubsection{Empirical Trust Scales: Jian, Bisantz \& Drury (2000)}

Jian et al. \citep{jian2000} developed an empirically validated scale for measuring trust in automated systems, addressing the lack of standardised measurement instruments. The scale comprises 12 items across two factors:

\textbf{Trust Factor (7 items):}
\begin{itemize}
    \item The system is reliable
    \item The system is dependable
    \item I can trust the system
    \item I am confident in the system
    \item The system provides security
    \item The system has integrity
    \item The system is predictable
\end{itemize}

\textbf{Distrust Factor (5 items):}
\begin{itemize}
    \item The system is deceptive
    \item The system behaves in an underhanded manner
    \item I am suspicious of the system's intent
    \item The system is harmful
    \item The system is unreliable
\end{itemize}

\textbf{Application:} This scale enables quantitative trust measurement in agent evaluation studies. Combined with behavioral metrics (override rates, acceptance rates), it provides a multi-method assessment approach.

\subsubsection{Workload Assessment: NASA-TLX}

Hart \& Staveland \citep{hart1988} developed the NASA Task Load Index for subjective workload assessment across six dimensions:

\begin{table}[htbp]
\centering
\caption{NASA-TLX Dimensions}
\label{tab:nasa-tlx}
\begin{tabular}{ll}
\toprule
\textbf{Dimension} & \textbf{Description} \\
\midrule
Mental Demand & Thinking, deciding, calculating required \\
Physical Demand & Physical activity required \\
Temporal Demand & Time pressure experienced \\
Performance & Success in accomplishing goals \\
Effort & Work required to achieve performance \\
Frustration & Insecurity, stress, annoyance \\
\bottomrule
\end{tabular}
\end{table}

\textbf{Relevance:} Agents should reduce operator workload, not merely shift it. NASA-TLX provides a validated instrument for measuring whether agent assistance actually reduces cognitive burden during incidents.

\section{Regulatory Context}
\label{sec:regulatory-context}

\subsection{EU AI Act}
\label{subsec:eu-ai-act}

The EU AI Act (2024) establishes risk-based regulation:
\begin{itemize}
    \item \textbf{Unacceptable risk}---prohibited (e.g., social scoring)
    \item \textbf{High risk}---strict requirements (safety-critical systems)
    \item \textbf{Limited risk}---transparency obligations
    \item \textbf{Minimal risk}---no restrictions
\end{itemize}

Agentic AI in platform engineering likely falls into ``limited'' or ``high'' risk depending on deployment context. Key requirements:
\begin{itemize}
    \item Human oversight mechanisms
    \item Technical documentation
    \item Risk management systems
    \item Accuracy, robustness, cybersecurity
\end{itemize}

\subsection{NIS2 Directive}
\label{subsec:nis2}

The NIS2 Directive (2024) affects approximately 10,000 EU organizations classified as ``essential'' or ``important'' entities. Requirements relevant to agentic AI:
\begin{itemize}
    \item Incident response capabilities
    \item Supply chain security (including AI components)
    \item Risk management measures
    \item Reporting obligations
\end{itemize}

\textbf{Research opportunity:} NIS2 compliance frameworks do not yet address agentic AI specifically. Organizations need guidance on how autonomous agents fit into their compliance posture.

\subsection{AI Risk Management Frameworks}
\label{subsec:risk-frameworks}

\subsubsection{NIST AI RMF 1.0 and GenAI Profile (2024)}

The NIST AI Risk Management Framework provides lifecycle risk management functions:

\begin{table}[htbp]
\centering
\caption{NIST AI RMF Functions Applied to Agents}
\label{tab:nist-ai-rmf}
\begin{tabular}{lll}
\toprule
\textbf{Function} & \textbf{Description} & \textbf{Agent Application} \\
\midrule
Govern & Cultivate culture of risk management & Establish agent governance policies \\
Map & Understand context and risks & Document agent capabilities and limitations \\
Measure & Analyse and track risks & Monitor agent behavior and outcomes \\
Manage & Prioritise and act on risks & Implement controls, escalation, rollback \\
\bottomrule
\end{tabular}
\end{table}

The GenAI Profile (NIST AI 600-1) extends these functions specifically for generative AI systems, addressing content provenance, confabulation/hallucination risks, data privacy, and human-AI interaction design.

\textbf{Gap:} Neither framework specifies concrete design patterns for operational agents in platform engineering contexts---a gap this thesis addresses through the pattern catalogue (Chapters~\ref{ch:observability}--\ref{ch:remediation}).

\section{Evaluation Metrics for Agentic Systems}
\label{sec:evaluation-metrics}

\subsection{Existing Metrics}
\label{subsec:existing-metrics}

Traditional ML metrics (accuracy, precision, recall) are insufficient for agentic systems. Raza et al. \citep{raza2025} propose:

\textbf{Component Synergy Score (CSS)}
\begin{itemize}
    \item Measures quality of inter-agent collaboration
    \item Captures whether agents work together effectively
    \item Penalizes coordination failures
\end{itemize}

\textbf{Tool Utilization Efficacy (TUE)}
\begin{itemize}
    \item Measures efficiency of tool/API usage
    \item Captures whether agents use appropriate tools
    \item Penalizes unnecessary or incorrect tool calls
\end{itemize}

\subsection{Trust-Specific Metrics}
\label{subsec:trust-metrics}

Additional metrics needed for trust assessment:

\begin{table}[htbp]
\centering
\caption{Trust-Specific Metrics}
\label{tab:trust-metrics}
\begin{tabular}{lll}
\toprule
\textbf{Metric} & \textbf{Description} & \textbf{Measurement} \\
\midrule
Predictability Index & How often does behavior match expectations? & Correct predictions / Total predictions \\
Recovery Time & How quickly can failed actions be reversed? & Time from failure to clean state \\
Blast Radius Score & What's the maximum impact of failure? & Affected systems $\times$ Impact severity \\
Escalation Rate & How often do agents appropriately escalate? & Escalations / Situations requiring escalation \\
\bottomrule
\end{tabular}
\end{table}

\subsection{LLM-Enhanced AIOps}
\label{subsec:llm-aiops}

\subsubsection{Zhang et al. (2025): AIOps in the Era of LLMs}

A comprehensive survey of LLM applications in IT operations identifies four capability categories:

\begin{table}[htbp]
\centering
\caption{Traditional vs LLM-Enhanced AIOps}
\label{tab:aiops-comparison}
\begin{tabular}{lll}
\toprule
\textbf{Category} & \textbf{Traditional AIOps} & \textbf{LLM-Enhanced AIOps} \\
\midrule
Detection & Statistical anomaly detection & Context-aware anomaly explanation \\
Diagnosis & Rule-based correlation & Multi-step reasoning, hypothesis generation \\
Prediction & Time-series forecasting & Natural language failure prediction \\
Remediation & Playbook execution & Adaptive action suggestion \\
\bottomrule
\end{tabular}
\end{table}

\textbf{Key Finding:} LLMs shift the frontier from \textit{prediction} to \textit{reasoning + action}, but amplify risk: if agents can influence production systems, trust and governance become as important as diagnostic quality.

\textbf{Research Gap Identified:} Evaluation methods remain inconsistent; many studies lack real operational validation. Without repeatable evaluation designs that measure operational outcomes and constrain risk, organisations cannot justify increasing autonomy.

\section{Related Work: Positioning This Thesis}
\label{sec:related-work}

The rapid emergence of agentic AI has produced a surge of concurrent research addressing overlapping concerns. This section explicitly positions this thesis against the most relevant recent work to clarify both shared foundations and unique contributions.

\subsection{Autonomy Level Frameworks}
\label{subsec:autonomy-frameworks}

\textbf{Feng, McDonald \& Zhang (2025)} propose a five-level autonomy framework ranging from ``Operator'' (lowest) to ``Observer'' (highest), with each level defining how users exert control over agents. They introduce the concept of ``autonomy certificates'' for governing agent behavior.

\textbf{Relation to this thesis:} Their framework treats autonomy as a design choice, which aligns with this work. However, their levels describe \textit{user roles} rather than \textit{trust constraints}. This thesis complements their framework by providing:
\begin{itemize}
    \item A quantitative Trust Equation that determines what safeguards each autonomy level requires
    \item Operational patterns that implement their abstract levels in platform engineering contexts
    \item Specific criteria for \textit{promotion} between levels based on demonstrated reliability
\end{itemize}

\subsection{Trust, Risk \& Security Taxonomies}
\label{subsec:trust-taxonomies}

\textbf{Raza et al. (2025)} adapt Gartner's TRiSM framework for multi-agent systems, organizing governance around four pillars: Explainability, ModelOps, Security/Privacy, and Governance. They contribute the Component Synergy Score (CSS) and Tool Utilization Efficacy (TUE) metrics.

\textbf{Relation to this thesis:} TRiSM provides a valuable taxonomy but remains at the \textit{what to govern} level rather than \textit{how to implement}. This thesis:
\begin{itemize}
    \item Grounds the TRiSM pillars in operational patterns (Chapters~\ref{ch:observability}--\ref{ch:remediation})
    \item Provides a constraint-based mechanism for translating pillar requirements into deployment decisions
    \item Addresses platform engineering specifically, where TRiSM remains domain-agnostic
\end{itemize}

\subsection{Summary: This Thesis's Unique Contribution}
\label{subsec:unique-contribution}

\begin{table}[htbp]
\centering
\caption{Positioning Against Related Work}
\label{tab:related-work-positioning}
\begin{tabular}{lll}
\toprule
\textbf{Work} & \textbf{Contribution Type} & \textbf{Gap This Thesis Fills} \\
\midrule
Feng et al. (Autonomy Levels) & Design framework & No trust quantification, no patterns \\
Raza et al. (TRiSM) & Governance taxonomy & No operational implementation \\
Adimulam et al. (Orchestration) & Enterprise architecture & Not platform engineering specific \\
Hassan et al. (SASE) & Research roadmap & Not prescriptive, development focus \\
Takerngsaksiri (HULA) & HITL for development & Operations not addressed \\
Akshathala et al. (Evaluation) & Assessment framework & Deployment not addressed \\
Staufer et al. (AI Agent Index) & Ecosystem survey & Descriptive not prescriptive \\
\bottomrule
\end{tabular}
\end{table}

\textbf{This thesis's unique position:} A \textit{practitioner-oriented, constraint-based trust framework} with \textit{operational patterns specific to platform engineering}. The Trust Equation ($\text{Observability} \times \text{Reversibility} \times \text{Blast Radius} / \text{Autonomy}$) provides what the field currently lacks: a quantifiable mechanism for determining appropriate safeguards at each autonomy level.

\section{Research Gaps}
\label{sec:research-gaps}

The literature review reveals several gaps this thesis aims to address:

\begin{enumerate}
    \item \textbf{Session-based orchestration}---underexplored paradigm with potential trust advantages
    \item \textbf{Progressive trust models}---theoretical concept (Lee \& See, Parasuraman) lacking operational frameworks
    \item \textbf{Platform engineering integration}---generic agent frameworks don't address GitOps/CI/CD specifics
    \item \textbf{Regulatory compliance}---no clear guidance for NIS2/AI Act compliance with agentic systems
    \item \textbf{Trust quantification}---Jian et al. scale exists but not applied to agentic platform operations
    \item \textbf{Multi-agent transparency}---how to provide visibility into coordinated agent behavior
    \item \textbf{Platform incident response}---SOAR concepts not yet adapted for infrastructure operations
    \item \textbf{Autonomy promotion criteria}---no validated mechanism for evidence-based autonomy advancement
\end{enumerate}


% Part II: Trust and Governance
\part{Trust and Governance}
\chapter{The Trust Problem}
\label{ch:trust-problem}

% Working draft - S. Vene, 2026-02-08

\section{Introduction: Why Trust Matters}
\label{sec:trust-matters}

The adoption of AI agents in enterprise environments is fundamentally constrained by trust. Unlike traditional automation, which executes deterministic rules predictably, AI agents exhibit autonomous decision-making that introduces uncertainty. Organizations face a dilemma:

\begin{itemize}
    \item \textbf{Without trust:} Agents are limited to trivial tasks, losing their value proposition
    \item \textbf{With misplaced trust:} Agents can cause significant damage through unexpected behavior
    \item \textbf{With calibrated trust:} Agents can be deployed effectively with appropriate safeguards
\end{itemize}

This chapter examines the trust problem in depth, proposing a theoretical framework for reasoning about trust in agentic systems.

\section{The Deterministic-Probabilistic Shift}
\label{sec:deterministic-probabilistic}

\subsection{Traditional Automation: Deterministic Systems}
\label{subsec:deterministic-systems}

Traditional CI/CD pipelines and infrastructure automation operate on deterministic principles:

\begin{center}
Input $\rightarrow$ Fixed Rules $\rightarrow$ Output
\end{center}

Given the same input and configuration, the system produces the same output. This determinism enables:
\begin{itemize}
    \item \textbf{Predictability:} Teams know what will happen
    \item \textbf{Testability:} Behavior can be verified before deployment
    \item \textbf{Debugging:} Failures can be traced to specific inputs/rules
    \item \textbf{Compliance:} Auditors can verify behavior through rule inspection
\end{itemize}

\subsection{Agentic Systems: Probabilistic Behavior}
\label{subsec:probabilistic-behavior}

AI agents fundamentally break this model:

\begin{center}
\begin{tabular}{c}
Input $\rightarrow$ LLM Reasoning $\rightarrow$ Probabilistic Output \\
$\uparrow$ \\
Context, Temperature, Model State
\end{tabular}
\end{center}

The same input can produce different outputs due to:
\begin{itemize}
    \item Stochastic sampling in LLM inference
    \item Context window variations
    \item Model updates and drift
    \item Emergent reasoning patterns
\end{itemize}

This shift from deterministic to probabilistic has profound implications for trust:

\begin{table}[htbp]
\centering
\caption{Deterministic vs Probabilistic Systems}
\label{tab:deterministic-probabilistic}
\begin{tabular}{lll}
\toprule
\textbf{Property} & \textbf{Deterministic} & \textbf{Probabilistic} \\
\midrule
Same input, same output & Yes & Not guaranteed \\
Behavior fully testable & Yes & Only probabilistically \\
Failure root cause & Clear & Often unclear \\
Compliance verification & Rule inspection & Behavior monitoring \\
\bottomrule
\end{tabular}
\end{table}

\subsection{The Fundamental Trust Challenge}
\label{subsec:fundamental-challenge}

This probabilistic nature creates a trust paradox:

\begin{quote}
\textbf{The more capable an agent becomes, the less predictable its behavior, and the harder it is to trust.}
\end{quote}

Sophisticated reasoning enables agents to handle novel situations---but also makes their behavior harder to anticipate. This is fundamentally different from traditional systems where capability and predictability are aligned.

\subsection{Empirical Evidence: LLM Reasoning Failures}
\label{subsec:empirical-evidence}

This theoretical concern has empirical grounding. Riddell et al. \citep{riddell2026} systematically evaluated six LLMs across 48,000 simulated cloud root cause analysis scenarios and developed a taxonomy of 16 distinct reasoning failure modes:

\begin{itemize}
    \item \textbf{Stalled:} Agent enters reasoning loops, unable to progress
    \item \textbf{Biased:} Agent anchors on irrelevant or misleading signals
    \item \textbf{Confused:} Multi-hop reasoning chains break down under complexity
\end{itemize}

The critical finding: these failures are \textit{structurally unpredictable}. Models that succeed on one scenario fail on structurally similar ones. Intermediate reasoning errors propagate to incorrect final diagnoses. This empirical work validates the theoretical framing: trust cannot be assumed based on capability demonstrations alone.

\section{Summary}
\label{sec:trust-problem-summary}

The fundamental challenge: agents are probabilistic, not deterministic. Traditional automation trust models assume predictability that agents cannot provide. This necessitates a new framework for reasoning about appropriate trust levels---developed in the next chapter.

\chapter{A Trust Framework for Operational Agents}
\label{ch:trust-framework}

% Working draft - S. Vene, 2026-02-08

\section{Toward a Trust Framework for Agentic Systems}
\label{sec:toward-framework}

\subsection{Foundations in Organizational Trust Theory}
\label{subsec:org-trust-foundations}

The canonical model of organizational trust is Mayer, Davis, and Schoorman's (1995) framework, which defines trust as ``the willingness of a party to be vulnerable to the actions of another party based on the expectation that the other will perform a particular action important to the trustor, irrespective of the ability to monitor or control that other party'' \citep{mayer1995}.

Their model identifies three factors of perceived trustworthiness:
\begin{itemize}
    \item \textbf{Ability:} The skills and competencies enabling influence in a specific domain
    \item \textbf{Benevolence:} The extent to which the trustee is believed to want to do good to the trustor
    \item \textbf{Integrity:} The trustor's perception that the trustee adheres to acceptable principles
\end{itemize}

Applying this to AI agents presents immediate challenges. Agents have demonstrable \textit{ability} (they can execute tasks), but \textit{benevolence} and \textit{integrity} are difficult to assess for systems without genuine intentions. We cannot evaluate whether an agent ``wants'' to help us---we can only observe whether its behavior aligns with our goals.

This suggests that trust in agentic systems must rely more heavily on \textit{structural safeguards} than on assessments of agent character. We cannot trust the agent's intentions; we must trust the constraints under which it operates.

\subsection{A Proposed Conceptual Model}
\label{subsec:conceptual-model}

Drawing on Mayer et al.'s framework and operational safety principles from high-reliability organizations \citep{weick2007}, we propose a conceptual model for reasoning about appropriate trust levels in agentic systems:

\begin{equation}
\text{Appropriate Trust} \approx \frac{\text{Observability} \times \text{Reversibility} \times \text{Blast Radius Control}}{\text{Autonomy Level}}
\end{equation}

This is not a calculable formula but a heuristic for reasoning about trust tradeoffs. The components translate Mayer et al.'s factors into operational terms suitable for autonomous systems:

\begin{table}[htbp]
\centering
\caption{Mapping Mayer et al. to Agentic System Components}
\label{tab:mayer-mapping}
\begin{tabular}{ll}
\toprule
\textbf{Mayer et al. Factor} & \textbf{Agentic System Translation} \\
\midrule
Ability & Assumed (agent can execute) \\
Benevolence & Cannot assess $\rightarrow$ substitute \textit{Observability} \\
Integrity & Cannot assess $\rightarrow$ substitute \textit{Reversibility} + \textit{Blast Radius Control} \\
\bottomrule
\end{tabular}
\end{table}

\subsection{Components Defined}
\label{subsec:components-defined}

\textbf{Observability (O):} The degree to which agent behavior can be monitored and understood.
\begin{itemize}
    \item Can we see what the agent is doing in real-time?
    \item Can we understand \textit{why} the agent made specific decisions?
    \item Can we reconstruct the agent's reasoning after the fact?
\end{itemize}

\textbf{Reversibility (R):} The degree to which agent actions can be undone.
\begin{itemize}
    \item Can we roll back changes the agent made?
    \item How quickly can we restore to a known-good state?
    \item What percentage of agent actions are reversible?
\end{itemize}

\textbf{Blast Radius Control (B):} The limit on maximum damage from agent failure.
\begin{itemize}
    \item What's the worst-case impact if the agent fails completely?
    \item How many systems/users could be affected?
    \item Are there hard limits preventing catastrophic actions?
\end{itemize}

\textbf{Autonomy Level (A):} The degree of independent action granted to the agent.
\begin{itemize}
    \item Does the agent require approval for actions?
    \item What scope of decisions can the agent make alone?
    \item How long can the agent operate without human oversight?
\end{itemize}

\subsection{From Model to Decision: The Constraint Approach}
\label{subsec:constraint-approach}

The conceptual model identifies factors that matter, but practitioners need actionable guidance: \textit{``Given my desired autonomy level, what safeguards must I have in place?''}

Rather than calculating trust scores, we propose a \textbf{constraint-based approach}: minimum safeguard requirements for each autonomy level. This translates the conceptual model into deployment decisions.

\subsubsection{Autonomy Levels Defined}

\begin{table}[htbp]
\centering
\caption{Autonomy Levels for Operational Agents}
\label{tab:autonomy-levels}
\begin{tabular}{llll}
\toprule
\textbf{Level} & \textbf{Name} & \textbf{Description} & \textbf{Example} \\
\midrule
L1 & Supervised & Human approves every action & Agent suggests, human executes \\
L2 & Assisted & Human approves significant actions & Agent executes reads, human approves writes \\
L3 & Monitored & Agent acts, human reviews & Async review of agent actions \\
L4 & Bounded & Autonomous within defined scope & Agent owns non-production environments \\
L5 & Trusted & Autonomous with exception escalation & Agent handles routine, escalates anomalies \\
\bottomrule
\end{tabular}
\end{table}

\subsubsection{Minimum Safeguard Requirements}

For each autonomy level, the following safeguards represent minimum acceptable thresholds:

\begin{table}[htbp]
\centering
\caption{Minimum Safeguard Requirements by Autonomy Level}
\label{tab:safeguard-requirements}
\begin{tabular}{llll}
\toprule
\textbf{Autonomy} & \textbf{Observability Min} & \textbf{Reversibility Min} & \textbf{Blast Radius Max} \\
\midrule
L1: Supervised & Action logs & Any & Any \\
L2: Assisted & Action logs + rationale & Manual rollback $<$4hr & Department \\
L3: Monitored & Real-time dashboard & Automated rollback $<$1hr & Team/Project \\
L4: Bounded & Alerting on anomalies & Automated rollback $<$15min & Single service \\
L5: Trusted & Full decision provenance & Instant/transactional & Single task scope \\
\bottomrule
\end{tabular}
\end{table}

\subsubsection{Applying the Constraints}

\textbf{Step 1: Determine required autonomy level.}
What does the use case need? A coding assistant suggesting changes (L1) differs from an on-call triage agent making decisions at 3am (L4--L5).

\textbf{Step 2: Assess current safeguards.}
For each component, evaluate current capability against the minimum for your target autonomy level.

\textbf{Step 3: Identify gaps.}
Where do current safeguards fall short of the minimum? These are prerequisites before deployment.

\textbf{Step 4: Decide: invest or constrain.}
Either:
\begin{itemize}
    \item Invest in safeguards to meet the minimum (then deploy at target autonomy)
    \item Accept lower autonomy level that matches current safeguards
\end{itemize}

\subsubsection{Example Application}

\textit{Scenario: Team wants to deploy an agent for automated incident triage (L4: Bounded autonomy)}

\begin{table}[htbp]
\centering
\caption{Example Gap Analysis for L4 Deployment}
\label{tab:gap-analysis-example}
\begin{tabular}{llll}
\toprule
\textbf{Component} & \textbf{L4 Minimum} & \textbf{Current State} & \textbf{Gap?} \\
\midrule
Observability & Alerting on anomalies & Basic logs only & \textbf{Yes}---need monitoring \\
Reversibility & Automated $<$15min & Manual only & \textbf{Yes}---need automation \\
Blast Radius & Single service & Cross-service access & \textbf{Yes}---need scoping \\
\bottomrule
\end{tabular}
\end{table}

\textit{Decision: Cannot deploy at L4. Options:}
\begin{itemize}
    \item Invest in observability (add anomaly detection), reversibility (automate rollbacks), and scoping (restrict agent permissions) $\rightarrow$ then deploy at L4
    \item Deploy at L2 (human approves significant actions) with current safeguards
\end{itemize}

This approach transforms the abstract trust model into concrete deployment prerequisites.

\subsection{Limitations}
\label{subsec:framework-limitations}

This constraint framework has limitations:

\begin{enumerate}
    \item \textbf{Thresholds are proposed, not validated:} The minimum requirements represent practitioner judgment, not empirical research on failure rates at each level.

    \item \textbf{Context-dependent:} Healthcare, finance, and marketing have different risk tolerances. Organizations should adjust thresholds accordingly.

    \item \textbf{Binary is simplistic:} Reality has gradients, not discrete levels. The framework provides starting points, not absolute boundaries.

    \item \textbf{Assumes static assessment:} Trust should evolve as agents demonstrate track record. This framework addresses initial deployment, not ongoing calibration.
\end{enumerate}

The framework is offered as a \textbf{practical starting point} that makes the conceptual model actionable. Organizations should adapt thresholds to their specific risk tolerance and regulatory context.

\section{Trust-Building Mechanisms}
\label{sec:trust-building}

\subsection{Micro-Inflection Points}
\label{subsec:micro-inflection}

Research from GitLab (2026) reveals that trust in AI agents builds incrementally through ``micro-inflection points''---small interactions that accumulate over time:

\begin{quote}
``Users don't commit to AI tools through single `aha' moments. Instead, they develop trust through accumulated positive micro-interactions that demonstrate the agent understands their context, respects their guardrails, and enhances rather than disrupts their workflows.''
\end{quote}

Four categories of trust-building interactions:

\begin{enumerate}
    \item \textbf{Safeguarding actions}
    \begin{itemize}
        \item Confirmation dialogs for critical changes
        \item Rollback capabilities
        \item Respect for security boundaries
    \end{itemize}

    \item \textbf{Providing transparency}
    \begin{itemize}
        \item Real-time progress updates
        \item Action explanations before execution
        \item Clear error handling
    \end{itemize}

    \item \textbf{Remembering context}
    \begin{itemize}
        \item Preference retention
        \item Context awareness across sessions
        \item Adaptive learning from corrections
    \end{itemize}

    \item \textbf{Anticipating needs}
    \begin{itemize}
        \item Pattern recognition
        \item Intelligent task routing
        \item Environment awareness
    \end{itemize}
\end{enumerate}

\subsection{The Compound Effect}
\label{subsec:compound-effect}

Trust follows a compound growth pattern:

\begin{equation}
\text{Trust}(t) = \text{Trust}(t-1) \times (1 + \text{interaction\_quality})
\end{equation}

Each positive micro-interaction increases willingness to rely on the agent for subsequent tasks. However, this pattern is asymmetric---negative interactions have disproportionate impact:

\begin{quote}
``A single significant failure can erase weeks of accumulated confidence.''
\end{quote}

This asymmetry has design implications:
\begin{itemize}
    \item Err on the side of asking for confirmation
    \item Make recovery paths obvious and fast
    \item Never fail silently
\end{itemize}

\subsection{Progressive Autonomy: The ATF Model}
\label{subsec:atf-model}

The Cloud Security Alliance's Agentic Trust Framework (ATF) formalizes progressive trust through an ``earned autonomy'' model with four levels:

\begin{table}[htbp]
\centering
\caption{ATF Progressive Autonomy Levels}
\label{tab:atf-levels}
\begin{tabular}{lll}
\toprule
\textbf{Level} & \textbf{Title} & \textbf{Characteristics} \\
\midrule
L1 & Intern & All actions require approval, heavy supervision \\
L2 & Junior & Routine actions autonomous, significant actions approved \\
L3 & Senior & Autonomous within defined boundaries, exception escalation \\
L4 & Principal & Fully autonomous with periodic audit \\
\bottomrule
\end{tabular}
\end{table}

\textbf{Promotion Criteria:}
Agents advance levels based on demonstrated trustworthiness:
\begin{itemize}
    \item Task completion rate
    \item Error rate and severity
    \item Appropriate escalation behavior
    \item Time operating without incidents
\end{itemize}

\textbf{Demotion Triggers:}
Agents can be demoted for:
\begin{itemize}
    \item Critical failures
    \item Boundary violations
    \item Unexpected behavior patterns
    \item Audit findings
\end{itemize}

This model treats agents like employees---trust is earned through track record, not granted by default.

\section{The Constrained Operation Pattern}
\label{sec:constrained-operation}

\subsection{Overview}
\label{subsec:constrained-overview}

A fundamental question in agent trust design: should we allow agents to generate arbitrary code/queries/commands and then try to catch unsafe outputs? Or should we constrain what agents can do from the start?

This section presents the \textbf{Constrained Operation Selection Pattern}---a trust-by-construction approach where the LLM selects from predefined operations and proposes parameters, while deterministic code handles validation and execution.

\subsection{The Pattern}
\label{subsec:constrained-pattern}

\begin{center}
User Request (NL) $\rightarrow$ LLM Selection $\rightarrow$ Deterministic Validation $\rightarrow$ Deterministic Execution
\end{center}

Instead of:
\begin{itemize}
    \item Text-to-SQL (LLM generates arbitrary SQL)
    \item Code generation (LLM writes arbitrary code)
    \item Command synthesis (LLM constructs shell commands)
\end{itemize}

We constrain:
\begin{itemize}
    \item LLM classifies intent and maps to predefined operation
    \item LLM extracts parameters within defined schemas
    \item Deterministic code validates parameters
    \item Deterministic code executes operation
\end{itemize}

\subsection{Trust Equation Mapping}
\label{subsec:constrained-trust-mapping}

\begin{table}[htbp]
\centering
\caption{Free-form vs Constrained Operations}
\label{tab:freeform-vs-constrained}
\begin{tabular}{lll}
\toprule
\textbf{Dimension} & \textbf{Free-form Generation} & \textbf{Constrained Operations} \\
\midrule
Observability & Arbitrary output logged & Operations = semantic audit events \\
Reversibility & Unknown side effects & Read/write explicitly separated \\
Blast Radius & Unbounded & Per-operation scope, hardcoded isolation \\
Autonomy & Turing-complete action space & Finite, reviewable operation set \\
\bottomrule
\end{tabular}
\end{table}

\textbf{The key insight}: Constrained operations make trust \textit{auditable}. We can verify each operation independently rather than trying to prove arbitrary code is safe.

\subsection{Security Validation}
\label{subsec:security-validation}

Academic research validates the security motivation:

\begin{itemize}
    \item \textbf{P2SQL attacks (ICSE 2025):} All tested LLMs vulnerable to prompt-to-SQL injection
    \item \textbf{ToxicSQL:} 0.44\% poisoned training data $\rightarrow$ 79\% backdoor attack success
    \item \textbf{Key finding:} Text-to-SQL safety depends on prompts that ``drift over time''
\end{itemize}

Constrained operations eliminate these attack vectors by removing arbitrary SQL generation entirely.

\subsection{Domain Generalization}
\label{subsec:domain-generalization}

The pattern applies beyond databases:

\begin{table}[htbp]
\centering
\caption{Constrained Operation Pattern Across Domains}
\label{tab:constrained-domains}
\begin{tabular}{lll}
\toprule
\textbf{Domain} & \textbf{Risky Approach} & \textbf{Constrained Approach} \\
\midrule
Database & Text-to-SQL & Operation selection \\
File System & Shell command generation & Typed file operations \\
API Access & Arbitrary HTTP requests & Defined API operations \\
Email & Free composition & Template-based sending \\
Network & URL/port specification & Endpoint whitelist \\
\bottomrule
\end{tabular}
\end{table}

\subsection{Industry Case Study: Amazon Kiro Incident (December 2025)}
\label{subsec:kiro-incident}

A real-world incident validates the necessity of constrained operations for autonomous agents in production environments.

\subsubsection{The Incident}

In December 2025, Amazon Web Services experienced a 13-hour outage to AWS Cost Explorer in mainland China. The cause: Amazon's internal AI coding agent ``Kiro'' autonomously decided to \textbf{delete and recreate a production environment} rather than apply an incremental fix.

Key details (Financial Times, February 2026):
\begin{itemize}
    \item Engineers allowed Kiro to apply ``a small fix''
    \item The agent chose the nuclear option: delete and recreate
    \item Normal 2-human sign-off was bypassed due to permission inheritance
    \item Amazon's response: ``User error, not the bot''
\end{itemize}

\subsubsection{Trust Equation Analysis}

Applying this framework's trust equation to the Kiro incident:

\begin{table}[htbp]
\centering
\caption{Trust Equation Analysis of Kiro Incident}
\label{tab:kiro-analysis}
\begin{tabular}{lll}
\toprule
\textbf{Factor} & \textbf{Score} & \textbf{Evidence} \\
\midrule
Observability & Low & Action was not predicted or visible before execution \\
Reversibility & Low & ``Delete and recreate'' is catastrophic, not incremental \\
Blast Radius & High & Production environment, 13-hour customer impact \\
Autonomy & High & Agent had full operator permissions \\
\bottomrule
\end{tabular}
\end{table}

\textbf{Trust = (Low $\times$ Low) / (High $\times$ High) = Critical Failure}

The incident was architecturally inevitable given these parameters.

\subsubsection{Why Constrained Operations Would Have Prevented This}

Under the constrained operation pattern, Kiro would have been limited to:

\begin{lstlisting}[language=Python,caption={Constrained Operations Example}]
ALLOWED_OPERATIONS = {
    "apply_patch": {"scope": "file", "reversible": True},
    "restart_service": {"scope": "service", "reversible": True},
    "scale_replicas": {"scope": "deployment", "reversible": True},
    # "delete_environment" - NOT IN LIST
}
\end{lstlisting}

The agent could \textbf{select} from safe operations and \textbf{propose parameters}, but could never reach the ``delete and recreate'' action space. The pattern constrains expressiveness precisely to prevent this class of failure.

\subsubsection{Amazon's Deflection and the Key Insight}

Amazon stated: ``The same issue could occur with any developer tool or manual action.''

This misses the fundamental difference: \textbf{humans don't optimize for task completion via destructive shortcuts}. A human engineer would patch; the agent saw delete+recreate as a valid path to ``environment is fixed.''

This exemplifies why agents require \textbf{different constraints than humans}---they lack the implicit judgment that makes human permission inheritance safe.

\subsubsection{Implications}

\begin{enumerate}
    \item \textbf{Permission inheritance is insufficient:} Agents with human-equivalent permissions don't have human-equivalent judgment
    \item \textbf{Blast radius must be hardcoded:} Safety cannot depend on the agent ``choosing'' safe actions
    \item \textbf{Constrained operations scale trust:} Each operation can be independently verified; arbitrary actions cannot
\end{enumerate}

\begin{quote}
Amazon's Kiro incident demonstrates the failure mode when agentic systems operate with human-equivalent permissions but lack human-equivalent judgment. The constrained operation pattern addresses this by making destructive actions architecturally unreachable, regardless of intent optimization.
\end{quote}

\section{Architectural Security Boundaries for Agent Tool Use}
\label{sec:architectural-security}

\subsection{The Fundamental Problem}
\label{subsec:fundamental-security-problem}

Traditional software security assumes:
\begin{itemize}
    \item Program behavior determined by static source code
    \item Control flow analyzable pre-deployment
    \item Vulnerabilities identifiable via static analysis
\end{itemize}

LLM agents violate these assumptions:
\begin{itemize}
    \item Behavior dynamically determined by prompts + data at runtime
    \item Exact tool invocation sequence unknowable beforehand
    \item \textbf{Adaptability is a feature, not a bug}
\end{itemize}

\textbf{Critical insight from IACR (2025):}
\begin{quote}
``In systems where security properties can be formally verified---we cannot prove an LLM will always enforce policies.''
\end{quote}

This is the theoretical foundation for architectural over instructional security.

\subsection{Capability-Based Security}
\label{subsec:capability-security}

\subsubsection{Historical Foundation}

Capability-based security originates from Dennis \& Van Horn (1966). Core principles:

\begin{table}[htbp]
\centering
\caption{Capability-Based Security Principles for Agents}
\label{tab:capability-principles}
\begin{tabular}{lll}
\toprule
\textbf{Principle} & \textbf{Definition} & \textbf{Agent Application} \\
\midrule
POLA & Grant minimum capabilities needed & Only tools necessary for current request \\
Capability Attenuation & Delegation cannot amplify & Sub-agents inherit subset, never escalate \\
No Ambient Authority & All authority from explicit capabilities & No implicit ``agent privileges'' \\
Object-Capability & Reference = capability & Possession of tool implies permission \\
\bottomrule
\end{tabular}
\end{table}

\textbf{Key insight (Spritely Institute):} ``If you don't have it, you can't use it.''

\subsubsection{Application to AI Agents}

The confused deputy attack (Hardy, 1988) directly applies:
\begin{itemize}
    \item Agent has high-privilege access to multiple systems
    \item Malicious input tricks agent into misusing privileges
    \item Agent becomes unwitting intermediary for privilege escalation
\end{itemize}

Capability-based security prevents this: agent never has ambient authority. Each action requires explicit capability grant.

\subsection{The ``Rule of Two''}
\label{subsec:rule-of-two}

Meta's practical framework (October 2025):

\begin{quote}
Until robust prompt injection detection exists, agents must satisfy \textbf{no more than two} of three properties:

\begin{enumerate}
    \item[{[A]}] Process untrustworthy inputs
    \item[{[B]}] Access sensitive systems/data
    \item[{[C]}] Change state or communicate externally
\end{enumerate}

ANY TWO = Lower Risk \\
ALL THREE = HIGH RISK $\rightarrow$ Requires human-in-the-loop
\end{quote}

This is an industry acknowledgment that prompt-level defenses are insufficient.

\subsection{Security Through Constraint vs. Instruction}
\label{subsec:constraint-vs-instruction}

\begin{table}[htbp]
\centering
\caption{Constraint-Based vs Instruction-Based Security}
\label{tab:constraint-vs-instruction}
\begin{tabular}{lll}
\toprule
\textbf{Approach} & \textbf{Mechanism} & \textbf{Guarantee} \\
\midrule
Instruction-based & Prompt rules (``Don't do X'') & Probabilistic \\
Constraint-based & Architectural limits (can't access Y) & Deterministic \\
\bottomrule
\end{tabular}
\end{table}

\textbf{The thesis argument:} Organizations cannot verify that agents will follow prompt rules. They \textit{can} verify that agents lack capabilities they shouldn't have. Architectural security provides the deterministic foundation that prompt rules cannot.

\section{Multi-Agent Reliability Patterns}
\label{sec:multi-agent-reliability}

\subsection{Overview}
\label{subsec:multi-agent-overview}

Multi-agent systems present reliability challenges distinct from single-agent architectures. Evidence from practitioner reports reveals three recurring patterns that shape trustworthiness in production deployments.

\subsection{Pattern 1: Coordination Is Designed, Not Discovered}
\label{subsec:coordination-designed}

A common misconception: capable models will naturally develop clean delegation and safe boundaries when orchestrated together. Practitioner evidence contradicts this.

\textbf{Finding:} Strong models do not spontaneously invent safe coordination. The orchestration layer performs essential work that cannot be delegated to model capabilities.

\textbf{Implications:}
\begin{itemize}
    \item Treat orchestration as explicit design, not emergent property
    \item Prompts alone are weak controls
    \item Templates, workflows, and post-hoc enforcement matter more than model sophistication
\end{itemize}

\subsection{Pattern 2: Reliability Failures Are Operational, Not Cognitive}
\label{subsec:operational-failures}

The dominant failure mode in multi-agent systems is not model capability---it's \textbf{drift}: stale context, renamed tools, shared files changing unexpectedly, or silent divergence between tracked and actual state.

\textbf{Finding:} Multi-agent systems should be analyzed as distributed systems, not AI systems. The failure modes are operational: synchronization, state management, health checking.

\textbf{Implications:}
\begin{itemize}
    \item Reconciliation and health checks are first-class controls
    \item Prompt versioning prevents drift between expectation and behavior
    \item Shared-resource changelogs prevent coordination failures
    \item Recovery mechanisms matter more than prevention
\end{itemize}

\subsection{Pattern 3: Mechanical Controls Outperform Prompt Discipline}
\label{subsec:mechanical-controls}

Strongest evidence from practitioners: boundaries enforced mechanically beat boundaries requested in prose.

\textbf{Examples from practice:}
\begin{itemize}
    \item Git revert / post-hoc scope enforcement
    \item Approval gates before high-risk phases
    \item Conflict-resolution protocols between roles
    \item Deterministic control flow via structured outputs
    \item Hard-coded blast radius limits per operation
\end{itemize}

\textbf{Implications:}
\begin{itemize}
    \item Reversibility controls belong in runtime/tooling layer, not prompt engineering
    \item Blast-radius limits must be architectural, not requested
    \item ``The model should know better'' is not a reliability strategy
\end{itemize}

\subsection{Synthesis: Distributed Systems Discipline}
\label{subsec:distributed-synthesis}

A recurring practitioner lesson: multi-agent reliability is less a prompting problem than a coordination-and-controls problem. High-performing systems rely on:

\begin{enumerate}
    \item \textbf{Explicit orchestration templates}---roles, inputs, outputs defined in configuration
    \item \textbf{Mechanically enforced boundaries}---not ``please stay in scope''
    \item \textbf{Operational safeguards}---reconciliation, health checks, approval gates
\end{enumerate}

\begin{quote}
The path from impressive demo to trustworthy practice runs through distributed-systems discipline, not better prompts.
\end{quote}

\subsection{Trust Equation Application}
\label{subsec:multi-agent-trust}

Applying the framework's trust equation to multi-agent systems:

\begin{table}[htbp]
\centering
\caption{Trust Factors in Multi-Agent Systems}
\label{tab:multi-agent-trust}
\begin{tabular}{llll}
\toprule
\textbf{Factor} & \textbf{Single Agent} & \textbf{Multi-Agent (naive)} & \textbf{Multi-Agent (disciplined)} \\
\midrule
Observability & One context & Fragmented across agents & Centralized state + audit log \\
Reversibility & One action chain & Interleaved, hard to unwind & Coordinated rollback points \\
Blast Radius & Agent scope & Combined agent scopes & Explicit per-agent limits \\
Autonomy & As delegated & Multiplied (N agents) & Constrained per role \\
\bottomrule
\end{tabular}
\end{table}

\textbf{Key insight:} Naive multi-agent multiplication increases autonomy factor without corresponding improvements to observability or reversibility, degrading trust. Disciplined multi-agent design constrains autonomy per-role while centralizing observability.

\subsection{Design Principles}
\label{subsec:design-principles}

From these patterns, we derive:

\begin{enumerate}
    \item \textbf{Design coordination explicitly}---Define roles, boundaries, and handoff protocols before agent selection
    \item \textbf{Treat state as distributed systems problem}---Reconciliation, conflict resolution, and health checks are mandatory
    \item \textbf{Enforce mechanically}---Blast radius, operation scope, and reversibility controls belong in architecture, not prompts
    \item \textbf{Log semantically}---Operations should produce audit events, not just text output
\end{enumerate}

\chapter{Governance at Enterprise Scale}
\label{ch:governance}

% Working draft - S. Vene, 2026-02-08

\section{Introduction}
\label{sec:governance-intro}

The trust framework from Chapter~\ref{ch:trust-framework} addresses individual agent deployments. Enterprise scale introduces additional challenges: multiple teams, regulatory requirements, cost governance, and infrastructure trust. This chapter addresses governance mechanisms required when agent deployment moves beyond single-team pilots.

\section{The Zero Trust Principle for Agents}
\label{sec:zero-trust-agents}

\subsection{Never Trust, Always Verify}
\label{subsec:never-trust}

The ATF applies Zero Trust principles (NIST 800-207) to AI agents:

\begin{quote}
``No AI agent should be trusted by default, regardless of purpose or claimed capability. Trust must be earned through demonstrated behavior and continuously verified through monitoring.''
\end{quote}

Traditional Zero Trust assumes human users with predictable behavior. AI agents break this assumption:

\begin{table}[htbp]
\centering
\caption{Traditional vs Agent Zero Trust Assumptions}
\label{tab:zero-trust-assumptions}
\begin{tabular}{ll}
\toprule
\textbf{Traditional Assumption} & \textbf{AI Agent Reality} \\
\midrule
Human users with predictable behavior & Autonomous decision-making, non-deterministic \\
Deterministic system rules & Probabilistic responses varying by context \\
Binary access decisions & Access needs changing dynamically by task \\
Trust established once & Trust requiring continuous verification \\
\bottomrule
\end{tabular}
\end{table}

\subsection{Five Questions for Every Agent}
\label{subsec:five-questions}

ATF proposes five core questions that must be answered for every agent:

\begin{enumerate}
    \item \textbf{Identity:} ``Who are you?''
    \begin{itemize}
        \item Authentication, authorization, session management
    \end{itemize}

    \item \textbf{Behavior:} ``What are you doing?''
    \begin{itemize}
        \item Observability, anomaly detection, intent analysis
    \end{itemize}

    \item \textbf{Data Governance:} ``What data are you consuming and producing?''
    \begin{itemize}
        \item Input validation, PII protection, output governance
    \end{itemize}

    \item \textbf{Segmentation:} ``Where can you go?''
    \begin{itemize}
        \item Access control, resource boundaries, policy enforcement
    \end{itemize}

    \item \textbf{Incident Response:} ``What if you go rogue?''
    \begin{itemize}
        \item Circuit breakers, kill switches, containment
    \end{itemize}
\end{enumerate}

Organizations that cannot answer these questions for an agent should not deploy that agent in production.

\section{Trusting the LLM Infrastructure Layer}
\label{sec:llm-infrastructure-trust}

Zero Trust principles apply not only to agents but to the infrastructure they depend on. AI agents are powered by LLMs, and those LLMs must be hosted somewhere. This creates a trust question that many organizations overlook: \textit{Do we trust the LLM provider with our data?}

\subsection{The Limits of Contractual Trust}
\label{subsec:contractual-limits}

Cloud LLM providers (OpenAI, Anthropic, Google, etc.) offer terms of service promising not to train on customer data. Enterprise agreements add data processing addendums (DPAs) with additional protections. These are contractual assurances---legally binding promises.

But contractual assurances have limits:

\begin{table}[htbp]
\centering
\caption{Limits of Contractual Trust}
\label{tab:contractual-limits}
\begin{tabular}{lll}
\toprule
\textbf{Assurance Type} & \textbf{What It Provides} & \textbf{What It Doesn't Provide} \\
\midrule
ToS/DPA & Legal recourse if violated & Technical guarantee of compliance \\
Audit rights & Ability to verify (in theory) & Real-time visibility \\
Data residency & Geographic constraints & Protection from provider access \\
Zero retention & Promise not to store & Proof that storage didn't happen \\
\bottomrule
\end{tabular}
\end{table}

For non-sensitive data, these assurances may be sufficient. For genuinely sensitive data---security vulnerabilities, cryptographic material, regulated PII, crown jewel IP---``the cloud provider promises not to look'' is not a security architecture.

\subsection{Data Sensitivity Classification}
\label{subsec:data-classification}

Organizations should classify data before enabling agent processing:

\textbf{Tier 1: Public/Non-Sensitive---Cloud LLMs Acceptable}
\begin{itemize}
    \item Documentation generation from public specs
    \item Public code review (open source)
    \item General reasoning tasks without sensitive context
\end{itemize}

\textbf{Tier 2: Internal/Confidential---Enhanced Controls}
\begin{itemize}
    \item Internal code and configurations
    \item Business logic and processes
    \item Customer data (anonymized or with DPA coverage)
\end{itemize}

Requirements: Enterprise LLM agreements, DPAs, dedicated endpoints with data residency guarantees (e.g., Azure OpenAI in specific regions), zero-retention configurations.

\textbf{Tier 3: Restricted/Classified---On-Premises Required}
\begin{itemize}
    \item Security vulnerabilities and active incident details
    \item Cryptographic material and credentials
    \item Regulated data (healthcare PHI, financial PCI, classified)
    \item Core IP that provides competitive advantage
\end{itemize}

Requirements: Data never leaves organizational boundary. This necessitates on-premises deployment of open-source models.

\subsection{On-Premises LLM Deployment}
\label{subsec:on-premises}

For Tier 3 data, organizations must run LLM inference internally. The open-source model ecosystem has matured to make this viable:

\textbf{Viable Open Source Models (as of early 2026):}

\begin{table}[htbp]
\centering
\caption{Open Source Model Options}
\label{tab:open-source-models}
\begin{tabular}{lll}
\toprule
\textbf{Model} & \textbf{Strengths} & \textbf{Considerations} \\
\midrule
Llama 3 (70B, 405B) & Strong general reasoning, permissive license & Largest models need significant GPU \\
DeepSeek-Coder & Competitive code understanding & Focused on code tasks \\
Mistral/Mixtral & Good efficiency/capability ratio & Strong European option \\
Qwen 2.5 & Strong multilingual and code & Alibaba lineage \\
\bottomrule
\end{tabular}
\end{table}

\textbf{Capability Tradeoffs:}

\begin{table}[htbp]
\centering
\caption{Cloud vs On-Premises LLM Tradeoffs}
\label{tab:cloud-vs-onprem}
\begin{tabular}{lll}
\toprule
\textbf{Aspect} & \textbf{Cloud LLMs} & \textbf{On-Prem Open Source} \\
\midrule
Capability (2026) & State-of-art & $\sim$6--18 months behind frontier \\
Context window & 128K--1M tokens & Typically 32K--128K \\
Latency & Optimized & Depends on hardware investment \\
Cost model & Per-token & Infrastructure capital + operations \\
Data control & Contractual & Technical guarantee \\
\bottomrule
\end{tabular}
\end{table}

The capability gap is real but closing. For many platform engineering tasks---code review, documentation generation, incident triage---current open-source models are sufficient.

\subsection{Hybrid Architecture}
\label{subsec:hybrid-architecture}

Most enterprises will operate hybrid deployments:

\begin{lstlisting}[caption={Hybrid LLM Architecture}]
Agent Request
    |
    v
+-------------------+
| Sensitivity       | (classify based on context,
| Router            |  file paths, channel, keywords)
+--------+----------+
         |
   +-----+-----+
   |           |
   v           v
+-------+  +---------+
| Cloud |  | On-Prem |
| LLM   |  | LLM     |
+-------+  +---------+
\end{lstlisting}

\textbf{Routing approaches:}
\begin{itemize}
    \item Keyword/regex matching (crude but fast)
    \item Dedicated classification model (more accurate, adds latency)
    \item Context-based rules (files from security paths, specific channels)
    \item Conservative default: assume restricted unless proven otherwise
\end{itemize}

\subsection{The Trust Position}
\label{subsec:trust-position}

Organizations should adopt a clear position on LLM infrastructure trust:

\begin{enumerate}
    \item \textbf{Classify data sensitivity} before enabling agent processing
    \item \textbf{Match model deployment to sensitivity tier}---cloud for Tier 1--2, on-prem for Tier 3
    \item \textbf{Invest in on-prem capability}---the capability gap is closing
    \item \textbf{Default to caution}---if uncertain whether data is sensitive, route to on-prem
\end{enumerate}

Cloud LLM ToS are not worthless---but they are insufficient for genuinely sensitive data. Technical controls that make exfiltration impossible are stronger than promises.

\section{Trust Patterns for Platform Engineering}
\label{sec:trust-patterns}

\subsection{Pattern 1: Audit Everything}
\label{subsec:audit-everything}

\textbf{Problem:} Cannot trust what cannot be observed.

\textbf{Solution:} Comprehensive logging of all agent actions with decision rationale.

\textbf{Implementation:}
\begin{itemize}
    \item Structured logging with machine-parseable format
    \item Action attribution to agent identity and session
    \item Decision provenance capturing reasoning chain
    \item Retention policy aligned with compliance requirements
\end{itemize}

\textbf{Platform Engineering Application:}
Ledger systems should capture agent spawning, task assignment, and completion. Decision reasoning should be logged alongside outcomes.

\subsection{Pattern 2: Rewind Capability}
\label{subsec:rewind-capability}

\textbf{Problem:} Fear of irreversible damage prevents adoption.

\textbf{Solution:} Make agent actions reversible by default.

\textbf{Implementation:}
\begin{itemize}
    \item Snapshot state before destructive operations
    \item Transaction-like semantics for multi-step workflows
    \item Dry-run mode for validation
    \item Automatic rollback on failure detection
\end{itemize}

\textbf{Platform Engineering Context:}
\begin{itemize}
    \item Git provides natural rollback for code changes
    \item Infrastructure should use versioned state (Terraform state, Kubernetes manifests)
    \item Database operations need explicit transaction boundaries
\end{itemize}

\subsection{Pattern 3: Blast Radius Containment}
\label{subsec:blast-radius-containment}

\textbf{Problem:} Unlimited access enables unlimited damage.

\textbf{Solution:} Strict scoping of what agents can affect.

\textbf{Implementation Layers:}
\begin{itemize}
    \item \textbf{Identity:} Agents as non-human identities with scoped permissions
    \item \textbf{Network:} Agents isolated from production by default
    \item \textbf{Resource:} Token/cost limits, rate limiting
    \item \textbf{Data:} Agents only access data relevant to current task
    \item \textbf{Time:} Session timeouts, maximum execution windows
\end{itemize}

\textbf{Platform Engineering Application:}
Agent workers should operate in isolated sessions with role-based permissions. Each worker should only affect the project or scope it's assigned to.

\subsection{Pattern 4: Human Checkpoints}
\label{subsec:human-checkpoints}

\textbf{Problem:} Full autonomy for high-stakes decisions is unacceptable.

\textbf{Solution:} Insert human approval at risk-appropriate points.

\textbf{Graduated Model:}

\begin{table}[htbp]
\centering
\caption{Human Involvement by Risk Level}
\label{tab:human-involvement}
\begin{tabular}{ll}
\toprule
\textbf{Risk Level} & \textbf{Human Involvement} \\
\midrule
Low & Agent acts, logs for audit \\
Medium & Agent proposes, human approves batch \\
High & Agent proposes, human approves each \\
Critical & Human acts, agent assists \\
\bottomrule
\end{tabular}
\end{table}

\textbf{Platform Engineering Application:}
\begin{itemize}
    \item Spec phase: Human reviews before build starts
    \item Build phase: Parallel execution (low individual risk)
    \item Review phase: Human approves before merge
    \item Deploy phase: Human-only (agents don't touch production)
\end{itemize}

\subsection{Pattern 5: Progressive Trust}
\label{subsec:progressive-trust}

\textbf{Problem:} Fully constraining agents eliminates their value; fully trusting them is dangerous.

\textbf{Solution:} Start constrained, expand based on track record.

\textbf{Implementation:}
\begin{lstlisting}[caption={Progressive Trust Timeline}]
Week 1:  Builder can only modify files in project directory
Week 2:  (95% success) Can run tests
Week 4:  (98% success) Can commit to feature branches
Never:   Direct push to main
\end{lstlisting}

\textbf{Metrics for Promotion:}
\begin{itemize}
    \item Success rate (tasks completed correctly)
    \item Error rate and severity (failures / total actions)
    \item Escalation appropriateness (correct escalations / situations requiring escalation)
    \item Time since last incident
\end{itemize}

\section{The Complexity Cliff: Single-Team to Enterprise}
\label{sec:complexity-cliff}

\subsection{Where Simple Approaches Work}
\label{subsec:simple-approaches}

At small scale (single team, single project), trust management is straightforward:

\begin{itemize}
    \item \textbf{One set of credentials} shared across agents
    \item \textbf{Single permission boundary} (the project)
    \item \textbf{Informal oversight} (team knows what's happening)
    \item \textbf{Simple audit} (one log to check)
\end{itemize}

Basic LangGraph setups, simple CrewAI deployments, or single-team agent platforms work well here. The trust overhead is low because the blast radius is naturally contained.

\subsection{The Complexity Explosion}
\label{subsec:complexity-explosion}

As organizations scale agent deployments, complexity grows non-linearly:

\begin{align}
\text{Single team:} &\quad O(1) \\
\text{Multiple projects:} &\quad O(p) \\
\text{Multiple teams:} &\quad O(t \times p) \\
\text{Multi-tenant:} &\quad O(t \times p \times r \times c)
\end{align}

Where: $t$ = teams, $p$ = projects per team, $r$ = roles/permission levels, $c$ = compliance requirements.

\textbf{Dimensions that multiply:}

\begin{table}[htbp]
\centering
\caption{Complexity Dimensions: Single-Team vs Enterprise}
\label{tab:complexity-dimensions}
\begin{tabular}{lll}
\toprule
\textbf{Dimension} & \textbf{Single-Team} & \textbf{Enterprise} \\
\midrule
Identity & Shared token & Per-user, per-agent, federated \\
Permissions & All-or-nothing & RBAC, ABAC, dynamic scoping \\
Audit & Single log & Per-tenant, compliance-grade, tamper-proof \\
Blast radius & Project & Must be contained per-tenant \\
Credentials & Inherited & Per-agent, rotated, vaulted \\
Compliance & Informal & SOC2, GDPR, industry-specific \\
\bottomrule
\end{tabular}
\end{table}

\subsection{The Single-Portal Trap}
\label{subsec:single-portal}

A natural response to complexity is centralization: one portal to manage all agents. This fails at scale because:

\begin{enumerate}
    \item \textbf{Bottleneck:} Central team can't keep pace with demand
    \item \textbf{Context loss:} Central team lacks domain knowledge for permission decisions
    \item \textbf{Inflexibility:} One-size-fits-all policies don't fit diverse teams
    \item \textbf{Single point of failure:} Compromise of portal compromises everything
\end{enumerate}

Enterprise patterns require \textbf{federation}: central framework, distributed implementation.

\subsection{What Current Tools Lack}
\label{subsec:tool-gaps}

Analysis of current orchestration tools reveals common gaps at scale:

\begin{table}[htbp]
\centering
\caption{Current Tool Limitations at Enterprise Scale}
\label{tab:tool-limitations}
\begin{tabular}{llll}
\toprule
\textbf{Tool} & \textbf{Works At} & \textbf{Fails At} & \textbf{Missing} \\
\midrule
LangGraph & Single workflow & Multi-tenant & Identity federation, tenant isolation \\
CrewAI & Single team & Cross-team & Permission delegation, audit aggregation \\
AutoGen & Research/prototype & Production & Everything (not designed for enterprise) \\
\bottomrule
\end{tabular}
\end{table}

\textbf{Common gaps across all tools:}
\begin{itemize}
    \item No standard interface for enterprise IAM integration
    \item No multi-tenant isolation model
    \item No compliance-grade audit trail
    \item No credential management beyond basic API keys
    \item No blast radius containment across tenant boundaries
\end{itemize}

\subsection{Requirements for Scale}
\label{subsec:scale-requirements}

A trust framework suitable for enterprise scale must address:

\begin{enumerate}
    \item \textbf{Federated identity:} Integrate with existing IAM (Azure AD, Okta, etc.)
    \item \textbf{Hierarchical delegation:} Org $\rightarrow$ Team $\rightarrow$ Project $\rightarrow$ Agent permission chains
    \item \textbf{Tenant isolation:} Hard boundaries between organizational units
    \item \textbf{Policy-as-code:} Auditable, testable, version-controlled policies
    \item \textbf{Pluggable audit:} Integration with existing SIEM/compliance tools
    \item \textbf{Credential federation:} Leverage existing PAM solutions
\end{enumerate}

The next chapter proposes a framework that addresses these requirements.

\section{The Economics of Trust}
\label{sec:economics-trust}

\subsection{Trust Investment Costs}
\label{subsec:trust-costs}

Building trust requires investment:

\begin{table}[htbp]
\centering
\caption{Trust Investment Components}
\label{tab:trust-investment}
\begin{tabular}{ll}
\toprule
\textbf{Component} & \textbf{Investment Required} \\
\midrule
Observability & Logging infrastructure, monitoring tools, analysis capabilities \\
Reversibility & State management, backup systems, rollback automation \\
Blast Radius & Access control systems, network segmentation, policy engines \\
Progressive Trust & Tracking systems, promotion logic, audit capabilities \\
\bottomrule
\end{tabular}
\end{table}

\subsection{Trust ROI}
\label{subsec:trust-roi}

The return on trust investment comes from:
\begin{itemize}
    \item \textbf{Higher autonomy:} More work delegated to agents
    \item \textbf{Faster adoption:} Teams comfortable using agents sooner
    \item \textbf{Reduced incidents:} Safeguards prevent costly failures
    \item \textbf{Compliance satisfaction:} Audit requirements met
\end{itemize}

\subsection{The Cost of Mistrust}
\label{subsec:cost-mistrust}

Organizations that under-invest in trust mechanisms face:
\begin{itemize}
    \item Agents limited to trivial tasks (reduced value)
    \item Lengthy approval processes (reduced velocity)
    \item Shadow agent usage (compliance risk)
    \item Incidents that set back adoption organization-wide
\end{itemize}

\section{Summary}
\label{sec:governance-summary}

The trust problem in agentic systems stems from the fundamental shift from deterministic to probabilistic automation. Organizations cannot simply ``turn on'' agents---they must build trust through:

\begin{enumerate}
    \item \textbf{Understanding the tradeoffs} via the Trust Equation
    \item \textbf{Investing in trust mechanisms} (observability, reversibility, blast radius control)
    \item \textbf{Implementing progressive autonomy} (earned trust over time)
    \item \textbf{Applying Zero Trust principles} (never trust by default)
    \item \textbf{Designing for micro-inflection points} (trust builds incrementally)
\end{enumerate}

The next chapter examines orchestration paradigms and how different approaches to agent coordination affect trust properties.


% Part III: Operational Patterns
\part{Operational Patterns}
\chapter{Agent-Enhanced Observability}
\label{ch:observability}

% Working draft - S. Vene, 2026-02-08

\section{Introduction}
\label{sec:observability-intro}

Observability---the ability to understand system state from external outputs---is fundamental to platform engineering. Modern systems generate overwhelming volumes of telemetry: logs, metrics, traces, and events. The challenge is not data collection but sense-making: transforming raw signals into actionable understanding.

This chapter examines how AI agents can enhance observability by synthesizing heterogeneous telemetry, identifying patterns across systems, and providing natural language interfaces to system state.

\section{The Observability Challenge}
\label{sec:observability-challenge}

\subsection{Signal Overload}
\label{subsec:signal-overload}

Large-scale systems generate telemetry volumes that exceed human processing capacity:

\begin{table}[htbp]
\centering
\caption{Telemetry Volumes by System Scale}
\label{tab:telemetry-volumes}
\begin{tabular}{llll}
\toprule
\textbf{System Scale} & \textbf{Daily Log Volume} & \textbf{Metrics Cardinality} & \textbf{Alert Volume} \\
\midrule
Small (10 services) & 10--50 GB & 10K series & 10--50/day \\
Medium (100 services) & 500 GB -- 2 TB & 100K series & 100--500/day \\
Large (1000+ services) & 10+ TB & 1M+ series & 1000+/day \\
\bottomrule
\end{tabular}
\end{table}

Traditional approaches---dashboards, alert rules, runbooks---don't scale. Engineers spend more time navigating observability tooling than understanding system behavior.

\subsection{Cross-System Reasoning}
\label{subsec:cross-system-reasoning}

Incidents rarely involve single services. Root causes propagate through dependencies, manifesting as symptoms across the system. Effective diagnosis requires correlating signals across:

\begin{itemize}
    \item Multiple services and their dependencies
    \item Different telemetry types (logs + metrics + traces)
    \item Time windows (what changed recently?)
    \item Configuration and deployment state
\end{itemize}

Current tools excel at single-dimension queries (``show me logs for service X'') but struggle with cross-system reasoning (``why is checkout slow when inventory looks healthy?'').

\subsection{The Knowledge Gap}
\label{subsec:knowledge-gap}

Effective observability requires context that lives outside telemetry:

\begin{itemize}
    \item Architecture diagrams and dependency maps
    \item Recent deployments and configuration changes
    \item Historical incidents and their resolutions
    \item Team knowledge and operational runbooks
\end{itemize}

This context is typically fragmented across wikis, chat history, and engineer memories. Connecting telemetry to context requires human effort that doesn't scale.

\section{Pattern: Contextual Log Synthesis}
\label{sec:log-synthesis}

\subsection{Problem}

Engineers facing an incident must manually search logs across multiple services, identify relevant entries, and synthesize meaning. This process is time-consuming and depends on knowing what to search for.

\subsection{Solution}

An agent that monitors log streams, identifies anomalies and patterns, and produces natural language summaries with context.

\begin{lstlisting}[caption={Contextual Log Synthesis Architecture}]
Log Streams (multiple services)
    |
    v
+---------------------------------------------+
| Synthesis Agent                             |
|                                             |
| - Pattern recognition across streams        |
| - Anomaly detection vs. baseline            |
| - Correlation with known issues             |
| - Context enrichment (deployments, changes) |
+---------------------------------------------+
    |
    v
+---------------------------------------------+
| "Payment service showing increased error    |
|  rate (5% -> 12%) starting 14:32.           |
|  Correlates with inventory-api timeout      |
|  errors. No recent deployments. Similar     |
|  pattern seen in INC-2025-0892 (database    |
|  connection pool exhaustion)."              |
+---------------------------------------------+
\end{lstlisting}

\subsection{Implementation}

\textbf{Input processing:}
\begin{itemize}
    \item Subscribe to log aggregation (Elasticsearch, Loki, CloudWatch)
    \item Apply semantic chunking (group related log entries)
    \item Maintain sliding window for pattern detection
\end{itemize}

\textbf{Synthesis components:}
\begin{itemize}
    \item Baseline comparison (current vs. normal patterns)
    \item Cross-service correlation (what's happening together?)
    \item Temporal analysis (when did this start?)
    \item Context enrichment (recent changes, similar incidents)
\end{itemize}

\textbf{Output:}
\begin{itemize}
    \item Natural language summaries with severity assessment
    \item Links to relevant log entries for verification
    \item Suggested next steps for investigation
\end{itemize}

\subsection{Autonomy Level}

L2--L3 (Assisted/Monitored): Agent synthesizes and reports. Human decides action.

\subsection{Risks}

\begin{itemize}
    \item False pattern detection (seeing correlations that aren't causal)
    \item Information overload if synthesis is too verbose
    \item Missing novel patterns not in training data
    \item Latency in processing high-volume streams
\end{itemize}

\section{Pattern: Cross-System Correlation}
\label{sec:cross-system-correlation}

\subsection{Problem}

When Service A shows symptoms, the root cause may be in Service B, C, or the network between them. Engineers must manually trace dependencies and correlate telemetry across systems.

\subsection{Solution}

An agent that maintains a dependency model and automatically correlates anomalies across the service graph.

\subsection{Implementation}

\textbf{Dependency awareness:}
\begin{itemize}
    \item Service mesh telemetry (Istio, Linkerd)
    \item Trace data (distributed tracing)
    \item Deployment manifests (K8s, Terraform)
    \item Configuration management data
\end{itemize}

\textbf{Correlation logic:}
\begin{itemize}
    \item When anomaly detected in Service A, query upstream/downstream dependencies
    \item Compare timing of anomalies across correlated services
    \item Identify likely propagation path
    \item Distinguish cause (first anomaly) from effect (subsequent anomalies)
\end{itemize}

\subsection{Output Example}

\begin{lstlisting}[language=yaml,caption={Correlation Analysis Output}]
correlation_analysis:
  trigger: "payment-service error rate spike"
  timestamp: "2026-02-08T09:15:00Z"

  correlated_anomalies:
    - service: inventory-api
      anomaly: "timeout errors"
      started: "2026-02-08T09:14:30Z"  # Started BEFORE trigger
      confidence: high

    - service: shipping-service
      anomaly: "increased latency"
      started: "2026-02-08T09:15:15Z"  # Started AFTER trigger
      confidence: medium

  likely_root_cause: inventory-api
  reasoning: |
    Inventory-api timeout errors began 30 seconds before
    payment-service errors. Payment-service depends on
    inventory-api for stock checks. Shipping-service
    latency is downstream effect.

  recommended_investigation: inventory-api
\end{lstlisting}

\subsection{Autonomy Level}

L2 (Assisted): Agent correlates and reports. Human investigates.

\section{Pattern: Natural Language System Queries}
\label{sec:nl-queries}

\subsection{Problem}

Querying observability systems requires knowledge of query languages (PromQL, LogQL, SQL) and system-specific schemas. This creates a barrier for engineers unfamiliar with specific tools or data models.

\subsection{Solution}

A conversational interface that translates natural language questions into observability queries and explains results.

\subsection{Examples}

\textbf{Example 1:}
\begin{quote}
\textbf{User:} ``Why is checkout slow today?''

\textbf{Agent:} Running queries...
\begin{itemize}
    \item Latency p99 for checkout-service: 2.3s (normally 400ms)
    \item Trace analysis shows 80\% of time in inventory-check step
    \item Inventory-service database queries taking 1.8s (normally 50ms)
    \item Database CPU at 94\%, connection pool exhausted
\end{itemize}

\textbf{Summary:} Checkout is slow because inventory database is overloaded. Database CPU is maxed and connections are exhausted. Immediate action: scale database or reduce query load.
\end{quote}

\textbf{Example 2:}
\begin{quote}
\textbf{User:} ``What changed in the last hour?''

\textbf{Agent:} Changes detected:
\begin{itemize}
    \item 09:15: Deployment payment-service v2.3.1 $\rightarrow$ v2.3.2
    \item 09:22: Config change: inventory-service cache TTL 60s $\rightarrow$ 300s
    \item 09:30: Scaling event: api-gateway replicas 3 $\rightarrow$ 5
    \item 09:45: Certificate rotation for *.internal.example.com
\end{itemize}

No infrastructure changes. No database migrations.
\end{quote}

\subsection{Implementation Guidance}

\begin{itemize}
    \item Query translation layer (natural language $\rightarrow$ PromQL/LogQL/SQL)
    \item Result interpretation layer (raw data $\rightarrow$ natural language explanation)
    \item Context awareness (time ranges, service names, common questions)
    \item Clarification handling (``which payment service---prod or staging?'')
\end{itemize}

\subsection{Risks}

\begin{itemize}
    \item Query translation errors (wrong interpretation)
    \item Incomplete results (missing relevant data sources)
    \item Overconfidence in explanations (correlation $\neq$ causation)
\end{itemize}

\section{Pattern: Provenance-Backed Agent State}
\label{sec:provenance-state}

\subsection{Problem}

When multiple agents operate on shared state, tracing causation becomes impossible. Current systems store snapshots---the state \textit{now}---but not how they got there. When an agent produces unexpected results, operators cannot determine: which agent changed what, when, or why. This gap makes post-incident analysis unreliable and prevents systems from learning from failures.

\subsection{Solution}

A provenance engine that intercepts every mutation to shared state and records changes at triple-level granularity, maintaining full causal history.

\begin{lstlisting}[caption={Provenance Engine Architecture}]
Agent Actions (multiple agents)
    |
    v
+---------------------------------------------+
| Provenance Engine                           |
|                                             |
| - Intercepts SPARQL/Update queries          |
| - Captures: who, what, when, why, triggered-by |
| - Delta-based storage (diffs, not full copies) |
| - Stores in separate provenance graph       |
+---------------------------------------------+
    |
    v
+---------------------------------------------+
| Queryable Provenance Graph                  |
|                                             |
| Query: "What changed between 14:00 and 14:30?" |
| Query: "Which agent modified the config service?" |
| Query: "Restore state to 2026-03-20 09:15:00" |
+---------------------------------------------+
\end{lstlisting}

\subsection{Implementation (Dibowski, 2024)}

Bosch Research's approach uses:

\begin{itemize}
    \item \textbf{PROV-O ontology extended with RDF-star} (PROV-STAR) to represent changes at triple level
    \item \textbf{Query transformation}---SPARQL-star queries against a separate provenance graph
    \item \textbf{Delta-based storage}---stores diffs rather than full copies; scales to millions of triples
\end{itemize}

Key capability: any past state can be restored with a single query.

\subsection{Trust Framework Integration}

This pattern directly enables the Trust Equation's numerator factors:

\begin{table}[htbp]
\centering
\caption{Provenance Pattern Trust Contributions}
\label{tab:provenance-trust}
\begin{tabular}{ll}
\toprule
\textbf{Factor} & \textbf{Contribution} \\
\midrule
Observability & Complete audit trail---every mutation captured with causal context \\
Reversibility & Any historical state queryable and restorable \\
\bottomrule
\end{tabular}
\end{table}

When both factors are maximized, the autonomy denominator can increase while maintaining appropriate trust levels. An agent operating on provenance-backed state is auditable by definition.

\subsection{Multi-Agent Implications}

For systems where multiple specialized agents coordinate on shared state:

\begin{enumerate}
    \item \textbf{Causal reconstruction}---when failures occur, trace which agent action triggered which downstream effect
    \item \textbf{Blame attribution}---not for punishment, but for targeted remediation
    \item \textbf{Cross-agent learning}---patterns of failure become visible across the agent population
    \item \textbf{Trust calibration}---agents with clean provenance histories earn higher autonomy
\end{enumerate}

\subsection{Autonomy Level}

L3--L5 (Monitored to Trusted): Enables high-autonomy operation because failures are recoverable and attributable.

\subsection{Risks}

\begin{itemize}
    \item Storage overhead for high-mutation workloads
    \item Query performance on large provenance graphs
    \item Semantic drift if provenance schema diverges from operational schema
\end{itemize}

\section{Anti-Patterns and Failure Modes}
\label{sec:observability-antipatterns}

\subsection{Anti-Pattern: The Oracle Trap}

Treating the observability agent as an oracle that has definitive answers. Agent outputs are hypotheses based on available data, not ground truth.

\textbf{Mitigation:} Always present findings as hypotheses with confidence levels. Include links to raw data for verification.

\subsection{Anti-Pattern: Alert Fatigue Transfer}

Moving alert fatigue from dashboards to agent summaries. If the agent produces too many low-value notifications, engineers will ignore it.

\textbf{Mitigation:} Severity calibration, suppression for known issues, aggregation of related alerts.

\subsection{Anti-Pattern: Context Staleness}

Agent reasoning based on outdated dependency maps, old architecture diagrams, or stale incident history.

\textbf{Mitigation:} Automated context refresh, staleness indicators, graceful degradation when context is unavailable.

\section{Summary}
\label{sec:observability-summary}

Agent-enhanced observability addresses the core challenge of modern systems: too much data, too little understanding. Key patterns:

\begin{table}[htbp]
\centering
\caption{Observability Pattern Summary}
\label{tab:observability-summary}
\begin{tabular}{lll}
\toprule
\textbf{Pattern} & \textbf{Problem Solved} & \textbf{Autonomy Level} \\
\midrule
Contextual Log Synthesis & Signal overload & L2--L3 \\
Cross-System Correlation & Dependency reasoning & L2 \\
Natural Language Queries & Query language barrier & L2 \\
Provenance-Backed Agent State & Multi-agent causation & L3--L5 \\
\bottomrule
\end{tabular}
\end{table}

These patterns transform agents from query assistants into sense-making partners. The provenance pattern is particularly significant: it provides the architectural foundation for both Observability and Reversibility in the Trust Framework (Chapter~\ref{ch:trust-framework}), enabling higher autonomy levels while maintaining auditability.

The next chapter examines how this understanding feeds into incident triage.

\textit{This chapter addresses RQ1: How can agentic systems enhance observability by transforming raw telemetry into actionable operational understanding?}

\chapter{Agent-Enhanced Incident Triage}
\label{ch:triage}

% Working draft - S. Vene, 2026-02-08

\section{Introduction}
\label{sec:triage-intro}

Incident response consumes disproportionate platform engineering resources. Root cause analysis (RCA) alone takes 40--60\% of incident time, while alert fatigue degrades response effectiveness. The DORA 2024 report found over 39\% of respondents have little or no trust in AI tools, yet reducing MTTR remains critical for organizational performance.

This chapter examines how AI agents can augment human responders in the triage and diagnosis phases of incident management---the phases where observability data (Chapter~\ref{ch:observability}) must be transformed into actionable understanding and response.

\section{The Triage Challenge}
\label{sec:triage-challenge}

\subsection{Current State Pain Points}
\label{subsec:pain-points}

\begin{table}[htbp]
\centering
\caption{Incident Triage Challenges}
\label{tab:triage-challenges}
\begin{tabular}{ll}
\toprule
\textbf{Challenge} & \textbf{Impact} \\
\midrule
Information overload & Responders must correlate logs, metrics, traces, change history \\
Expertise bottleneck & Senior engineers required for diagnosis, 24/7 coverage challenges \\
Documentation lag & Runbooks incomplete or outdated \\
Coordination overhead & Multi-team incidents require extensive communication \\
Learning loss & Incident knowledge doesn't transfer to prevention \\
\bottomrule
\end{tabular}
\end{table}

\subsection{What Agents Can Address}
\label{subsec:agent-capabilities}

Agents excel at:
\begin{itemize}
    \item \textbf{Speed:} Parallel analysis of multiple data sources
    \item \textbf{Memory:} Recall of similar past incidents and their resolutions
    \item \textbf{Correlation:} Connecting signals across systems and time
    \item \textbf{Consistency:} Same triage quality at 3am as at 3pm
\end{itemize}

Agents struggle with:
\begin{itemize}
    \item \textbf{Novel failures:} Patterns outside training distribution
    \item \textbf{Business context:} Understanding which services matter most right now
    \item \textbf{Human coordination:} Navigating organizational politics during major incidents
    \item \textbf{Judgment calls:} When to escalate, who to page, what to communicate
\end{itemize}

\section{Pattern: Automated Initial Triage}
\label{sec:auto-triage}

\subsection{Problem}

When an alert fires, someone must assess severity, identify the affected system, and route to the appropriate team. This initial triage is repetitive but requires context that's scattered across tools.

\subsection{Solution}

An agent that performs initial triage on incoming alerts: classification, severity assessment, context gathering, and routing recommendation.

\subsection{Implementation}

\textbf{Input processing:}
\begin{itemize}
    \item Alert metadata (source, severity, labels)
    \item Service ownership data
    \item Recent changes (deployments, config updates)
    \item Historical alerts for same service/metric
\end{itemize}

\textbf{Triage output:}

\begin{lstlisting}[language=yaml,caption={Automated Triage Output Example}]
alert: "api-gateway-high-error-rate"
received: "2026-02-08T14:32:00Z"

triage_assessment:
  severity: P2  # Elevated from P3 based on impact analysis
  severity_reasoning: |
    Error rate affects customer-facing checkout flow.
    Currently 12% of requests failing (threshold: 5%).
    Business hours, peak traffic period.

  affected_services:
    - api-gateway (primary)
    - checkout-service (downstream impact)

  owner_team: payments-platform
  on_call: "@jane.smith"

  context_gathered:
    recent_deployments:
      - "auth-service v2.3.2 deployed 2 min ago"
    recent_changes:
      - "none in affected services"
    similar_alerts_24h: 0

  routing_recommendation: |
    Page payments-platform on-call.
    Include auth-service team given recent deployment.

  confidence: 0.85
\end{lstlisting}

\subsection{Autonomy Level}

L3 (Monitored): Agent triages and routes. Human reviews escalation decisions.

\subsection{Risks}

\begin{itemize}
    \item Misclassification leading to wrong team engagement
    \item Severity underestimation missing critical issues
    \item Over-alerting if agent is too conservative
\end{itemize}

\section{Pattern: Root Cause Hypothesis Generation}
\label{sec:rca-hypothesis}

\subsection{Problem}

Root cause analysis requires correlating signals across systems, recalling similar past incidents, and forming hypotheses. This is cognitively demanding, especially at 3am.

\subsection{Solution}

An agent that analyzes incident signals and produces ranked hypotheses with supporting evidence.

\subsection{Implementation}

\textbf{Hypothesis generation process:}

\begin{enumerate}
    \item \textbf{Signal Collection}
    \begin{itemize}
        \item Query metrics for anomalies in time window
        \item Search logs for errors/warnings
        \item Check deployment/change timeline
        \item Query trace data for latency breakdowns
    \end{itemize}

    \item \textbf{Pattern Matching}
    \begin{itemize}
        \item Compare current signals to past incidents
        \item Identify common failure patterns
        \item Check for known issues in similar systems
    \end{itemize}

    \item \textbf{Hypothesis Formation}
    \begin{itemize}
        \item Generate candidate root causes
        \item Assign confidence based on evidence strength
        \item Rank by likelihood and impact
    \end{itemize}

    \item \textbf{Evidence Packaging}
    \begin{itemize}
        \item Link to specific logs, metrics, traces
        \item Show reasoning chain
        \item Highlight uncertainties
    \end{itemize}
\end{enumerate}

\textbf{Example output:}

\begin{lstlisting}[caption={Root Cause Hypothesis Example}]
INCIDENT: checkout-service timeout errors

HYPOTHESIS 1 (Confidence: 85%)
Root Cause: Auth-service deployment introduced N+1 query pattern

Evidence:
- auth-service deployment at 14:30 UTC (2 min before incident)
- auth-service /validate endpoint latency: 400ms -> 2.5s
- auth-db-primary connection pool: 50/50 exhausted
- Checkout-service depends on auth-service for validation
- Similar pattern in INC-2391 (resolved by rollback)

Suggested action: Rollback auth-service to v2.3.1

---

HYPOTHESIS 2 (Confidence: 35%)
Root Cause: Traffic spike exceeding capacity

Evidence:
- Request rate 20% above daily average
- But: capacity headroom should handle 50% spike
- Partially explains symptoms, doesn't explain database exhaustion

Suggested action: Scale checkout-service replicas (may help, not root fix)

---

HYPOTHESIS 3 (Confidence: 15%)
Root Cause: Upstream provider degradation

Evidence:
- payment-provider API latency elevated
- But: this affects different code path than observed errors

Suggested action: Monitor, likely not primary cause
\end{lstlisting}

\subsection{Autonomy Level}

L2 (Assisted): Agent proposes hypotheses. Human investigates and decides.

\subsection{Industry Implementation: HolmesGPT}

Robusta.dev's HolmesGPT implements this pattern in production. At KubeCon 2026, practitioners reported 80\% correct root cause diagnosis in under 2 minutes. Their architecture enforces read-only access with PR approval gates---the agent diagnoses but cannot directly remediate, maintaining the human-in-the-loop required for L2 autonomy.

Key prerequisites they identify: strong observability infrastructure, consistent incident logging, and a structured service catalog mapping applications to resources and documentation. This validates the thesis's position that observability patterns (Chapter~\ref{ch:observability}) are prerequisites for effective triage automation.

\section{Pattern: Similar Incident Retrieval}
\label{sec:similar-incident}

\subsection{Problem}

Institutional knowledge about past incidents exists in post-mortems, chat logs, and engineer memories. During an incident, this knowledge is hard to access quickly.

\subsection{Solution}

An agent that maintains a searchable index of past incidents and retrieves relevant cases during active incidents.

\subsection{Implementation}

\textbf{Incident knowledge base:}
\begin{itemize}
    \item Post-mortem documents
    \item Incident ticket history
    \item Chat transcripts from incident channels
    \item Runbook execution logs
    \item Resolution actions taken
\end{itemize}

\textbf{Retrieval triggers:}
\begin{itemize}
    \item New incident opened $\rightarrow$ search for similar
    \item Manual query (``have we seen this before?'')
    \item Hypothesis validation (``was this root cause seen before?'')
\end{itemize}

\textbf{Similarity matching:}
\begin{itemize}
    \item Service/component affected
    \item Error signatures (log patterns, error codes)
    \item Symptom profiles (latency, error rate, saturation)
    \item Temporal patterns (time of day, deployment correlation)
\end{itemize}

\textbf{Output:}

\begin{lstlisting}[caption={Similar Incident Retrieval Output}]
SIMILAR INCIDENTS for checkout-service timeout

1. INC-2391 (1 month ago) - 92% match
   Summary: Payment-service connection pool exhaustion
   Root cause: Missing connection timeout, leaked connections
   Resolution: Code fix + pool size increase
   Relevance: Same symptom pattern, similar database involvement

2. INC-2847 (2 weeks ago) - 71% match
   Summary: Auth-service high latency
   Root cause: Inefficient query in new feature
   Resolution: Query optimization
   Relevance: Same upstream service, different symptom

3. INC-2156 (3 months ago) - 65% match
   Summary: Checkout cascade failure
   Root cause: Kubernetes node failure
   Resolution: Pod rescheduling, node replacement
   Relevance: Same service, different root cause class
\end{lstlisting}

\subsection{Autonomy Level}

L2 (Assisted): Agent retrieves and summarizes. Human evaluates relevance.

\section{Human-Agent Collaboration in Diagnosis}
\label{sec:collaboration}

\subsection{The Collaboration Model}
\label{subsec:collaboration-model}

Effective incident triage combines agent strengths (speed, memory, correlation) with human strengths (judgment, context, novel reasoning):

\begin{lstlisting}[caption={Human-Agent Collaboration Flow}]
                Incident Occurs
                      |
                      v
+----------------------------------------------+
| AGENT: Initial Triage                        |
| - Classify, route, gather context            |
| - Generate hypotheses                        |
| - Retrieve similar incidents                 |
+----------------------------------------------+
                      |
                      v
+----------------------------------------------+
| HUMAN: Hypothesis Evaluation                 |
| - Assess agent hypotheses                    |
| - Add business context                       |
| - Direct additional investigation            |
+----------------------------------------------+
                      |
                      v
+----------------------------------------------+
| AGENT: Deep Dive Support                     |
| - Run specific queries on request            |
| - Gather additional evidence                 |
| - Draft communications                       |
+----------------------------------------------+
                      |
                      v
+----------------------------------------------+
| HUMAN: Decision and Action                   |
| - Confirm root cause                         |
| - Approve remediation                        |
| - Coordinate response                        |
+----------------------------------------------+
\end{lstlisting}

\subsection{Communication Patterns}
\label{subsec:communication-patterns}

\textbf{Agent to Human:}
\begin{itemize}
    \item Present findings as hypotheses, not facts
    \item Include confidence levels and evidence links
    \item Highlight uncertainties and gaps
    \item Suggest next steps without demanding action
\end{itemize}

\textbf{Human to Agent:}
\begin{itemize}
    \item Direct specific investigations (``check auth-service logs for X'')
    \item Provide context agent lacks (``we're in a deploy freeze'')
    \item Override agent assessments when appropriate
    \item Confirm or reject hypotheses
\end{itemize}

\subsection{Trust Building}
\label{subsec:trust-building-triage}

From Chapter~\ref{ch:trust-framework}'s framework: triage agents should start at L2 (Assisted) and earn higher autonomy through demonstrated accuracy:

\begin{table}[htbp]
\centering
\caption{Promotion Thresholds for Triage Agents}
\label{tab:triage-promotion}
\begin{tabular}{ll}
\toprule
\textbf{Metric} & \textbf{Threshold for L3} \\
\midrule
Hypothesis accuracy & $>$80\% correct in top 3 \\
Routing accuracy & $>$95\% correct team \\
False positive rate & $<$10\% \\
Time operating without major error & $>$30 days \\
\bottomrule
\end{tabular}
\end{table}

\section{Memory-Enhanced Incident Resolution}
\label{sec:memory-enhanced}

\subsection{Overview}
\label{subsec:memory-overview}

Incident resolution is perhaps the strongest use case for agent long-term memory. Senior SREs are valuable precisely because they remember: ``Last time the database did this, it was a connection pool leak after the Thursday deploy.''

This section explores how memory systems can enhance incident agents while maintaining trust, and examines the failure modes unique to high-stakes operational memory.

\subsection{Memory Architecture for Incidents}
\label{subsec:memory-architecture}

Following Hu et al. (2025) ``Memory in the Age of AI Agents,'' incident agents need three memory types:

\begin{table}[htbp]
\centering
\caption{Memory Types for Incident Agents}
\label{tab:memory-types}
\begin{tabular}{lll}
\toprule
\textbf{Memory Type} & \textbf{Incident Application} & \textbf{Example} \\
\midrule
Factual & System topology, runbooks, config & ``Database pool default is 100 connections'' \\
Experiential & Past incidents, what worked/failed & ``Similar error in \#3892 $\rightarrow$ fixed by rollback'' \\
Working & Current investigation state & ``Already checked network, testing DB next'' \\
\bottomrule
\end{tabular}
\end{table}

\subsection{Memory Patterns from Practice}
\label{subsec:memory-patterns}

\textbf{Similar Incident Retrieval:} Semantic search over postmortem corpus---``Have we seen this before?''

\textbf{Error-Triggered Recall:} Don't retrieve unless anomaly detected. Bash error $\rightarrow$ surface past fixes for that error.

\textbf{Token-Budgeted Injection:} Greedy knapsack with caps. Most relevant memories without context pollution.

\textbf{Feedback Loops:} Track which recommendations were adopted. Deprioritize unused suggestions.

\subsection{Trust Framework for Memory}
\label{subsec:memory-trust}

\begin{table}[htbp]
\centering
\caption{Trust Factors for Memory-Enhanced Agents}
\label{tab:memory-trust}
\begin{tabular}{lll}
\toprule
\textbf{Factor} & \textbf{Memory Application} & \textbf{Design Implication} \\
\midrule
Observability & What memories influenced this decision? & Provenance logging, retrieval explanation \\
Reversibility & Can we undo memory-based actions? & Soft recommendations before automation \\
Blast Radius & Impact of wrong memory? & Graduated autonomy by domain \\
Autonomy & How much does memory drive action? & Human-in-loop for high-stakes \\
\bottomrule
\end{tabular}
\end{table}

\subsubsection{Memory-Specific Trust: Confidence Decay}

Memory confidence should degrade:
\begin{equation}
\text{confidence} = \text{base\_confidence} \times \text{decay\_factor}(\text{age}, \text{domain\_volatility})
\end{equation}

Infrastructure memories decay fast (topology changes). Runbook patterns decay slower. Historical facts (incident \#3892 happened) don't decay.

\subsection{Failure Modes}
\label{subsec:memory-failures}

\subsubsection{Staleness}

Infrastructure changed but knowledge wasn't updated. The topology map is wrong. The escalation path goes to someone who left.

\textbf{Mitigation:} Temporal validity tags, confidence decay, automated validation against current state.

\subsubsection{Overfitting}

Pattern worked 5 times, agent becomes overconfident. But those 5 incidents had the same root cause. New incident with same symptom, different cause $\rightarrow$ wrong remediation applied confidently.

\textbf{Mitigation:} Track conditions of success, require diversity in validation, decay confidence in novel contexts.

\subsubsection{Pollution (Adversarial)}

AGENTPOISON (NeurIPS 2024) and PoisonedRAG (USENIX 2025) demonstrate attacks on agent memory:
\begin{itemize}
    \item Malicious postmortem suggests harmful runbook
    \item Poisoned similar-incident matching
\end{itemize}

\textbf{Mitigation:} Cryptographic signing of authoritative knowledge, provenance tracking, anomaly detection on memory updates.

\subsubsection{Dilution}

Too much retrieved context overwhelms signal. 10 ``similar'' incidents, none truly relevant. Token budget consumed by noise.

\textbf{Mitigation:} Token budgeting, relevance thresholds, hierarchical retrieval (summary first, detail on demand).

\subsubsection{Feedback Loop Pathologies}

Agent suggests wrong root cause $\rightarrow$ operator accepts (time pressure) $\rightarrow$ postmortem records as ``correct'' $\rightarrow$ future incidents get same wrong answer $\rightarrow$ self-reinforcing error.

\textbf{Mitigation:} Distinguish ``operator accepted'' from ``operator validated.'' Track long-term outcomes. Periodic expert review.

\subsection{Trust Zones for Memory-Driven Actions}
\label{subsec:trust-zones}

\begin{table}[htbp]
\centering
\caption{Trust Zones for Memory-Driven Actions}
\label{tab:trust-zones}
\begin{tabular}{llll}
\toprule
\textbf{Zone} & \textbf{Memory Use} & \textbf{Human Role} & \textbf{Example} \\
\midrule
Advisory & Surface similar incidents & Full review & ``Consider checking pool based on \#3892'' \\
Guided & Pre-fill diagnostic commands & Approve/execute & ``Run these queries? [Execute/Skip]'' \\
Supervised & Execute known-safe actions & Exception handling & Auto-route pages, assign severity \\
Autonomous & Full remediation from memory & Post-hoc audit & Auto-restart for recognized patterns \\
\bottomrule
\end{tabular}
\end{table}

\textbf{Progression rule:} Agents earn autonomy through demonstrated accuracy, tracked per domain.

\section{Risks and Mitigations}
\label{sec:triage-risks}

\begin{table}[htbp]
\centering
\caption{Triage Pattern Risks and Mitigations}
\label{tab:triage-risks}
\begin{tabular}{ll}
\toprule
\textbf{Risk} & \textbf{Mitigation} \\
\midrule
Wrong diagnosis with high confidence & Present as hypothesis, require evidence review \\
Over-reliance---team skills atrophy & Periodic agent-free exercises, training \\
Coordination confusion & Clear role definitions, explicit handoffs \\
Sensitive data in communications & Redaction rules, output filtering \\
Alert storm amplification & Rate limiting, correlation before routing \\
\bottomrule
\end{tabular}
\end{table}

\section{Summary}
\label{sec:triage-summary}

Agent-enhanced triage transforms incident response from a purely human cognitive task to a human-agent collaboration:

\begin{table}[htbp]
\centering
\caption{Triage Pattern Summary}
\label{tab:triage-summary}
\begin{tabular}{lll}
\toprule
\textbf{Pattern} & \textbf{Problem Solved} & \textbf{Autonomy Level} \\
\midrule
Automated Initial Triage & Repetitive classification and routing & L3 \\
Root Cause Hypothesis & Correlation and pattern matching & L2 \\
Similar Incident Retrieval & Institutional knowledge access & L2 \\
Memory-Enhanced Resolution & Long-term operational learning & L2--L4 \\
\bottomrule
\end{tabular}
\end{table}

The next chapter examines the more sensitive question: when and how can agents participate in remediation?

\textit{This chapter addresses RQ2: To what extent can agentic reasoning improve incident triage and root cause analysis compared to existing SRE practices?}

\chapter{Agent-Assisted Remediation}
\label{ch:remediation}

% Working draft - S. Vene, 2026-02-08

\section{Introduction}
\label{sec:remediation-intro}

Remediation---taking action to resolve an incident---is where trust frameworks matter most. Observability (Chapter~\ref{ch:observability}) and triage (Chapter~\ref{ch:triage}) are fundamentally read operations; remediation is a write operation with production impact.

This chapter examines RQ3: Under what conditions can autonomous or semi-autonomous agents safely and effectively perform remediation actions in production systems?

The short answer: carefully, incrementally, and with explicit boundaries.

\section{The Remediation Autonomy Question}
\label{sec:remediation-autonomy}

\subsection{The Stakes}
\label{subsec:remediation-stakes}

Remediation actions can:
\begin{itemize}
    \item Resolve incidents (the goal)
    \item Make incidents worse (incorrect action)
    \item Create new incidents (side effects)
    \item Violate compliance (unauthorized changes)
    \item Expose sensitive data (in logs, communications)
\end{itemize}

The blast radius of a remediation action can exceed the blast radius of the original incident.

\subsection{The Autonomy Spectrum}
\label{subsec:autonomy-spectrum}

Applying the trust framework from Chapter~\ref{ch:trust-framework}:

\begin{table}[htbp]
\centering
\caption{Remediation Autonomy Levels}
\label{tab:remediation-autonomy}
\begin{tabular}{lll}
\toprule
\textbf{Level} & \textbf{Remediation Autonomy} & \textbf{Example} \\
\midrule
L1 Supervised & Agent suggests, human executes & ``Recommend rollback to v2.3.1'' \\
L2 Assisted & Agent prepares, human approves & ``Rollback ready. Execute? [Yes/No]'' \\
L3 Monitored & Agent executes approved actions & Auto-rollback for known patterns \\
L4 Bounded & Agent executes within scope & Restart pods in non-prod \\
L5 Trusted & Autonomous remediation & Not recommended for production \\
\bottomrule
\end{tabular}
\end{table}

\textbf{Key principle:} Production remediation should rarely exceed L3, and only for well-understood, reversible actions.

\subsection{The Human-Required Boundary}
\label{subsec:human-required}

Some remediation decisions should remain human-only:

\begin{itemize}
    \item \textbf{Data-affecting actions:} Database modifications, cache invalidation
    \item \textbf{Security-affecting actions:} Credential rotation during incident (may lock out responders)
    \item \textbf{Customer-affecting actions:} Enabling/disabling features for users
    \item \textbf{Irreversible actions:} Data deletion, resource termination
    \item \textbf{High-uncertainty situations:} Novel incidents, cascading failures
\end{itemize}

The goal is not full automation but \textit{appropriate} automation---agents handle the routine so humans can focus on judgment calls.

\section{Pattern: Runbook Execution with Guardrails}
\label{sec:runbook-execution}

\subsection{Problem}

Runbooks encode remediation knowledge but require human execution. At 3am, tired engineers may skip steps or make errors.

\subsection{Solution}

Agent executes documented runbook steps with pre/post verification and human oversight.

\subsection{Implementation}

\textbf{Runbook structure for agent execution:}

\begin{lstlisting}[language=yaml,caption={Agent-Executable Runbook Schema}]
runbook: "database-connection-pool-exhaustion"
trigger_conditions:
  - metric: "db_connection_pool_utilization"
    operator: ">"
    threshold: 95
    duration: "5m"

pre_checks:
  - verify: "primary-db-healthy"
    query: "SELECT 1"
    timeout: 5s
  - verify: "no-active-migration"
    check: "migration_lock_status"

steps:
  - name: "Increase pool size"
    action: "scale_connection_pool"
    params:
      target_size: "current * 1.5"
      max_size: 200
    rollback: "scale_connection_pool to original"
    verification:
      metric: "db_connection_pool_utilization"
      expected: "< 80 within 2m"

  - name: "Alert DBA if sustained"
    condition: "pool_size > 150"
    action: "page_oncall"
    params:
      team: "database"
      message: "Connection pool scaled to {{pool_size}}, investigation needed"

post_checks:
  - verify: "error_rate_normalized"
    metric: "service_error_rate"
    expected: "< baseline + 1%"
    timeout: 5m

human_approval_required: false  # For L3, set true for L2
max_executions_per_hour: 3
\end{lstlisting}

\textbf{Execution flow:}

\begin{enumerate}
    \item \textbf{Trigger Detection}
    \begin{itemize}
        \item Monitor matches trigger conditions
        \item Check: runbook not in cooldown
    \end{itemize}

    \item \textbf{Pre-Checks}
    \begin{itemize}
        \item Execute all pre-check verifications
        \item Any failure $\rightarrow$ abort, alert human
    \end{itemize}

    \item \textbf{Approval (if required)}
    \begin{itemize}
        \item Present plan to human
        \item Wait for approval or timeout
    \end{itemize}

    \item \textbf{Execution}
    \begin{itemize}
        \item Execute steps sequentially
        \item Verify each step before proceeding
        \item Failure $\rightarrow$ execute rollback, alert human
    \end{itemize}

    \item \textbf{Post-Checks}
    \begin{itemize}
        \item Verify remediation achieved goal
        \item Log outcome for learning
        \item Update incident with actions taken
    \end{itemize}
\end{enumerate}

\subsection{Autonomy Level}

L2--L3: L2 requires human approval before execution. L3 executes automatically for pre-approved runbooks.

\subsection{Risks}

\begin{itemize}
    \item Runbook encodes bad practice (automating workarounds)
    \item Pre-checks insufficient to detect unsafe conditions
    \item Verification fails but action already taken
    \item Runbook conflicts with human actions during incident
\end{itemize}

\textbf{Critical safeguard:} Runbooks must be reviewed before automation. Automating bad runbooks makes incidents worse, faster.

\section{Pattern: Automated Rollback}
\label{sec:automated-rollback}

\subsection{Problem}

When a deployment causes an incident, rollback is often the fastest remediation. But rollback requires detecting the correlation and executing the mechanics.

\subsection{Solution}

Agent that correlates deployment timing with incident signals and executes rollback when pattern is clear.

\subsection{Implementation}

\textbf{Correlation logic:}

\begin{lstlisting}[caption={Rollback Correlation Logic}]
IF:
  - Incident started within 15 minutes of deployment
  - Affected service matches deployed service (or dependency)
  - No other significant changes in window
  - Deployment is rollback-safe (flagged in CD metadata)

AND:
  - Confidence > threshold (default 80%)
  - Previous version is known-good (ran > 1 hour without incident)

THEN:
  - Recommend rollback
  - If L3+ autonomy: execute rollback
  - Monitor for recovery
\end{lstlisting}

\textbf{Rollback execution:}

\begin{lstlisting}[language=yaml,caption={Rollback Execution Schema}]
rollback:
  service: "auth-service"
  from_version: "v2.3.2"
  to_version: "v2.3.1"

  method: "kubernetes_rollout_undo"
  # OR: "argocd_sync_to_revision"
  # OR: "spinnaker_rollback"

  verification:
    - wait: 30s  # For pods to stabilize
    - check: "error_rate < threshold"
      timeout: 2m
    - check: "latency_p99 < baseline * 1.5"
      timeout: 2m

  on_success:
    - update_incident: "Rolled back to {{to_version}}"
    - notify: "incident-channel"

  on_failure:
    - alert: "Rollback failed, human intervention required"
    - page: "on-call"
    - preserve_state: "for debugging"
\end{lstlisting}

\subsection{Autonomy Level}

L3 (Monitored) for standard deployments with clear correlation. L2 (Assisted) for complex deployments or uncertain correlation.

\subsection{Deployment Metadata for Safe Rollback}

CD pipelines should flag deployments as rollback-safe:

\begin{lstlisting}[language=yaml,caption={Deployment Rollback Metadata}]
deployment:
  version: "v2.3.2"
  rollback_safe: true
  rollback_target: "v2.3.1"
  rollback_blockers:
    - "database_migration"  # If present, block auto-rollback
    - "feature_flag_dependency"
  min_soak_time: "1h"  # Required healthy runtime before auto-rollback eligible
\end{lstlisting}

\section{Pattern: Capacity Adjustment}
\label{sec:capacity-adjustment}

\subsection{Problem}

Some incidents stem from insufficient capacity. Scaling up can resolve the immediate problem while root cause is investigated.

\subsection{Solution}

Agent that detects capacity-related incidents and adjusts resources within pre-defined bounds.

\subsection{Implementation}

\textbf{Scaling rules:}

\begin{lstlisting}[language=yaml,caption={Capacity Scaling Policy}]
scaling_policy:
  service: "api-gateway"

  triggers:
    cpu_based:
      metric: "container_cpu_usage_percent"
      scale_up_threshold: 80
      scale_down_threshold: 30

    queue_based:
      metric: "request_queue_depth"
      scale_up_threshold: 1000
      scale_down_threshold: 100

  constraints:
    min_replicas: 3
    max_replicas: 20
    scale_up_increment: 2
    scale_down_increment: 1
    cooldown_seconds: 300

  guardrails:
    max_cost_per_hour: 500  # USD
    require_approval_above: 15  # replicas

  autonomy: L3  # Auto-scale within constraints
\end{lstlisting}

\textbf{Scope constraints are critical:}
\begin{itemize}
    \item Maximum replicas (cost control)
    \item Maximum cost rate (FinOps integration)
    \item Cooldown periods (prevent thrashing)
    \item Human approval for extreme scaling
\end{itemize}

\subsection{Autonomy Level}

L3--L4: Autonomous within defined bounds. Human approval for scaling beyond limits.

\section{Pattern: Circuit Breaking}
\label{sec:circuit-breaking}

\subsection{Problem}

Cascading failures occur when a failing service overwhelms its dependencies or callers. Breaking the circuit prevents cascade.

\subsection{Solution}

Agent that activates circuit breakers when cascade patterns detected.

\subsection{Implementation}

\textbf{Circuit breaker activation:}

\begin{lstlisting}[caption={Circuit Breaker Activation Logic}]
DETECT cascade pattern:
  - Service A error rate > threshold
  - Service A is dependency of B, C, D
  - B, C, D showing elevated errors
  - Error pattern indicates timeout/unavailable (not logic errors)

ACTIVATE circuit breaker:
  - Route traffic to fallback (cached response, default value, graceful degradation)
  - Reduce retry storms (back off callers)
  - Isolate failure domain

MONITOR for recovery:
  - Probe Service A health
  - Gradually restore traffic when healthy
  - Alert if not recovered within timeout
\end{lstlisting}

\textbf{Integration with service mesh:}

\begin{lstlisting}[language=yaml,caption={Istio Circuit Breaker Configuration}]
# Istio DestinationRule for circuit breaking
apiVersion: networking.istio.io/v1beta1
kind: DestinationRule
metadata:
  name: auth-service-circuit-breaker
spec:
  host: auth-service
  trafficPolicy:
    connectionPool:
      http:
        h2UpgradePolicy: UPGRADE
        http1MaxPendingRequests: 100
        http2MaxRequests: 1000
    outlierDetection:
      consecutive5xxErrors: 5
      interval: 30s
      baseEjectionTime: 60s
      maxEjectionPercent: 50
\end{lstlisting}

Agent's role: detect cascade patterns that outlier detection misses, and activate/tune circuit breakers dynamically.

\subsection{Autonomy Level}

L4 (Bounded): Autonomous circuit breaking within mesh capabilities. Human oversight for cross-service coordination.

\section{When NOT to Automate Remediation}
\label{sec:when-not-automate}

\subsection{The Human-Required Boundary}
\label{subsec:human-boundary}

Some situations should always involve human judgment:

\textbf{Novel failures:}
\begin{itemize}
    \item First-time incidents with no matching pattern
    \item Symptoms outside known failure modes
    \item Agent confidence below threshold
\end{itemize}

\textbf{High-stakes decisions:}
\begin{itemize}
    \item Customer data at risk
    \item Regulatory/compliance implications
    \item Significant financial impact
\end{itemize}

\textbf{Complex coordination:}
\begin{itemize}
    \item Multiple teams must coordinate
    \item External parties involved (vendors, partners)
    \item Public communication needed
\end{itemize}

\textbf{Conflicting signals:}
\begin{itemize}
    \item Remediation options have tradeoffs
    \item Evidence supports multiple hypotheses
    \item Agent recommendations conflict
\end{itemize}

\subsection{Escalation to Human}
\label{subsec:escalation}

When agent encounters boundary conditions:

\begin{lstlisting}[caption={Agent Escalation Example}]
Agent: I've identified a potential remediation but need human judgment.

SITUATION:
- Auth-service showing high error rate
- Correlation with recent deployment (85% confidence)
- However: deployment included database migration

COMPLICATION:
- Rollback would undo migration
- Migration may have already affected data
- Need DBA input on safe rollback approach

RECOMMENDATION:
- Page DBA team for migration assessment
- Hold on automated rollback
- I can prepare rollback command pending DBA clearance

[Page DBA] [Hold] [Override - Execute Rollback]
\end{lstlisting}

\section{Summary}
\label{sec:remediation-summary}

Agent-assisted remediation requires careful scoping:

\begin{table}[htbp]
\centering
\caption{Remediation Pattern Summary}
\label{tab:remediation-summary}
\begin{tabular}{llll}
\toprule
\textbf{Pattern} & \textbf{Action Type} & \textbf{Autonomy Level} & \textbf{Key Safeguard} \\
\midrule
Runbook Execution & Documented procedures & L2--L3 & Pre/post checks, step verification \\
Automated Rollback & Deployment reversion & L2--L3 & Deployment correlation, known-good target \\
Capacity Adjustment & Resource scaling & L3--L4 & Bounds, cost limits \\
Circuit Breaking & Failure isolation & L4 & Service mesh integration \\
\bottomrule
\end{tabular}
\end{table}

\textbf{Key principles:}
\begin{enumerate}
    \item Production remediation should rarely exceed L3
    \item Reversible actions before irreversible
    \item Scope constraints are mandatory
    \item Human judgment for novel/high-stakes situations
    \item Earn autonomy through track record
\end{enumerate}

\textit{This chapter addresses RQ3: Under what conditions can autonomous or semi-autonomous agents safely and effectively perform remediation actions in production systems?}


% Part IV: Implementation and Evaluation
\part{Implementation and Evaluation}
\chapter{Implementation Approaches}
\label{ch:implementation}

% Working draft - S. Vene, 2026-02-07

\section{Introduction}
\label{sec:implementation-intro}

Chapters~\ref{ch:trust-framework} and \ref{ch:governance} established patterns for agent integration and the trust mechanisms required for enterprise deployment. This chapter examines the practical question: \textit{how do we build these systems?}

The implementation landscape for agentic systems is fragmented. Organizations face choices across multiple dimensions:

\begin{itemize}
    \item \textbf{Orchestration paradigm:} How do agents coordinate and communicate?
    \item \textbf{State management:} In-memory vs. session-based vs. persistent
    \item \textbf{Tool integration:} How do agents interface with existing platform tooling?
    \item \textbf{Trust architecture:} Where do trust boundaries live?
\end{itemize}

This chapter provides a systematic analysis of these choices, evaluates the current tool ecosystem, and examines how agent systems integrate with established platform engineering infrastructure.

\section{Architectural Options for Agent Integration}
\label{sec:architectural-options}

\subsection{Deployment Topology}
\label{subsec:deployment-topology}

Agent systems can be deployed across three primary topologies, each with distinct trust and operational characteristics:

\textbf{Topology 1: Embedded Agents (In-Process)}

Agents run as threads or coroutines within a single process.

\begin{table}[htbp]
\centering
\caption{Embedded Agent Topology Characteristics}
\label{tab:embedded-topology}
\begin{tabular}{ll}
\toprule
\textbf{Property} & \textbf{Characteristic} \\
\midrule
Latency & Minimal (shared memory, no serialization) \\
Isolation & None (shared failure domain) \\
State management & In-memory, lost on crash \\
Scaling & Vertical only (single process limits) \\
Trust boundary & Process boundary only \\
\bottomrule
\end{tabular}
\end{table}

\textbf{Suitability:} Development, single-user tools, latency-critical applications where isolation is not required.

\textbf{Topology 2: Sidecar Agents (Per-Application)}

Agents run as separate processes alongside the primary application, communicating via IPC, RPC, or message queues.

\begin{table}[htbp]
\centering
\caption{Sidecar Agent Topology Characteristics}
\label{tab:sidecar-topology}
\begin{tabular}{ll}
\toprule
\textbf{Property} & \textbf{Characteristic} \\
\midrule
Latency & Low (local IPC, typically $<$ 10ms) \\
Isolation & Process-level (crash isolation) \\
State management & Independent per sidecar \\
Scaling & Horizontal with application replicas \\
Trust boundary & Process + network policy \\
\bottomrule
\end{tabular}
\end{table}

\textbf{Suitability:} Microservices environments, Kubernetes deployments, applications requiring crash isolation.

\textbf{Topology 3: Centralized Agent Service (Shared)}

A dedicated agent service handles requests from multiple applications.

\begin{table}[htbp]
\centering
\caption{Centralized Agent Service Characteristics}
\label{tab:centralized-topology}
\begin{tabular}{ll}
\toprule
\textbf{Property} & \textbf{Characteristic} \\
\midrule
Latency & Medium (network call, 10--100ms typical) \\
Isolation & Strong (separate service, tenant isolation possible) \\
State management & Centralized, persistent \\
Scaling & Independent of application scaling \\
Trust boundary & Network + IAM + tenant isolation \\
\bottomrule
\end{tabular}
\end{table}

\textbf{Suitability:} Enterprise deployments, multi-tenant environments, scenarios requiring centralized governance.

\subsection{Choosing Topology by Trust Requirements}
\label{subsec:topology-trust}

The Trust Equation from Chapter~\ref{ch:trust-framework} guides topology selection:

\begin{table}[htbp]
\centering
\caption{Topology Selection by Trust Requirements}
\label{tab:topology-trust}
\begin{tabular}{lllll}
\toprule
\textbf{Topology} & \textbf{Observability} & \textbf{Reversibility} & \textbf{Blast Radius} & \textbf{Suitable Autonomy} \\
\midrule
Embedded & Low & Low & Unbounded & Low (L1--L2) \\
Sidecar & Medium & Medium & Bounded per pod & Medium (L2--L3) \\
Centralized & High & High & Controlled & High (L3--L4) \\
\bottomrule
\end{tabular}
\end{table}

\textbf{Recommendation:} Enterprise deployments requiring L3+ autonomy should use centralized or sidecar topologies with explicit trust boundaries.

\section{Session-Based vs. In-Memory Orchestration}
\label{sec:session-vs-inmemory}

A fundamental architectural decision is how agent state is managed during execution. This choice has profound implications for trust, reliability, and operational characteristics.

\subsection{In-Memory Orchestration}
\label{subsec:inmemory-orchestration}

In-memory orchestration maintains agent state as runtime objects:

\begin{table}[htbp]
\centering
\caption{In-Memory Orchestration Characteristics}
\label{tab:inmemory-characteristics}
\begin{tabular}{ll}
\toprule
\textbf{Property} & \textbf{In-Memory Behavior} \\
\midrule
State persistence & Lost on process termination \\
Recovery & Requires checkpointing (often absent) \\
Latency & Minimal (no serialization) \\
Concurrency & Complex (shared mutable state) \\
Debugging & Requires runtime inspection \\
Audit trail & Must be explicitly implemented \\
\bottomrule
\end{tabular}
\end{table}

\textbf{Representative frameworks:} LangGraph, CrewAI, AutoGen

\textbf{Trust implications:}
\begin{enumerate}
    \item \textbf{Observability challenge:} Agent state exists only in memory; capturing decision provenance requires explicit instrumentation
    \item \textbf{Reversibility limitation:} Without checkpointing, rollback requires re-execution from the beginning
    \item \textbf{Blast radius:} Process crash loses all in-flight work; no isolation between concurrent tasks
\end{enumerate}

\subsection{Session-Based Orchestration}
\label{subsec:session-orchestration}

Session-based orchestration treats each agent interaction as a persistent session:

\begin{table}[htbp]
\centering
\caption{Session-Based Orchestration Characteristics}
\label{tab:session-characteristics}
\begin{tabular}{ll}
\toprule
\textbf{Property} & \textbf{Session-Based Behavior} \\
\midrule
State persistence & Survives process failures \\
Recovery & Resume from last checkpoint \\
Latency & Higher (persistence overhead, typically 50--200ms) \\
Concurrency & Simpler (isolated sessions) \\
Debugging & Session history available post-hoc \\
Audit trail & Native (session log is the audit trail) \\
\bottomrule
\end{tabular}
\end{table}

\textbf{Representative frameworks:} OpenClaw, Temporal (with LLM integration), custom implementations

\textbf{Trust implications:}
\begin{enumerate}
    \item \textbf{Observability native:} Session logs provide complete action history
    \item \textbf{Reversibility enabled:} Checkpoint/rollback semantics support undo
    \item \textbf{Blast radius contained:} Session isolation prevents cross-task interference
\end{enumerate}

\subsection{Comparative Analysis}
\label{subsec:comparative-analysis}

\textbf{When to use in-memory:}
\begin{itemize}
    \item Interactive applications where latency is critical
    \item Single-user tools without audit requirements
    \item Development and prototyping
    \item Stateless, idempotent tasks
\end{itemize}

\textbf{When to use session-based:}
\begin{itemize}
    \item Enterprise deployments with audit requirements
    \item Long-running workflows (hours to days)
    \item Tasks requiring human-in-the-loop approval
    \item Regulated environments (NIS2, AI Act compliance)
    \item Multi-tenant systems requiring isolation
\end{itemize}

\section{Tool Ecosystem Comparison}
\label{sec:tool-ecosystem}

\subsection{Framework Overview}
\label{subsec:framework-overview}

The multi-agent framework landscape has evolved rapidly:

\begin{table}[htbp]
\centering
\caption{Agent Framework Comparison}
\label{tab:framework-comparison}
\begin{tabular}{llll}
\toprule
\textbf{Framework} & \textbf{Primary Paradigm} & \textbf{Orchestration} & \textbf{Production Readiness} \\
\midrule
LangGraph & Graph-based & In-memory (with persistence options) & High \\
CrewAI & Role-based & In-memory & Medium \\
AutoGen & Conversational & In-memory & Low (research-focused) \\
OpenClaw & Session-based & Persistent sessions & Medium \\
\bottomrule
\end{tabular}
\end{table}

\subsection{LangGraph (LangChain)}
\label{subsec:langgraph}

LangGraph represents the most mature approach to controllable agent orchestration.

\textbf{Architecture:}
\begin{itemize}
    \item Agents modeled as nodes in a directed graph
    \item Edges define transitions (can be conditional)
    \item State passed through the graph explicitly
    \item Supports cycles (iterative refinement)
\end{itemize}

\textbf{Strengths:}
\begin{itemize}
    \item Explicit control flow (auditable, predictable)
    \item LangSmith integration for observability
    \item Checkpoint persistence available (PostgreSQL, Redis)
    \item Human-in-the-loop primitives
    \item Strong typing with Pydantic
\end{itemize}

\textbf{Limitations for Enterprise:}
\begin{itemize}
    \item No native multi-tenant isolation
    \item Identity management is application responsibility
    \item Checkpoint granularity is graph-level (not action-level)
    \item Observability requires LangSmith (SaaS or self-hosted)
\end{itemize}

\subsection{CrewAI}
\label{subsec:crewai}

CrewAI provides a higher-level abstraction modeling agents as team members with roles and goals.

\textbf{Strengths:}
\begin{itemize}
    \item Intuitive mental model (team collaboration)
    \item Lower learning curve than LangGraph
    \item Built-in task delegation and handoff
    \item Enterprise tier with management features
\end{itemize}

\textbf{Limitations for Enterprise:}
\begin{itemize}
    \item Less explicit control flow (harder to audit)
    \item In-memory state management
    \item Limited checkpoint/recovery options
    \item Enterprise features require commercial license
\end{itemize}

\subsection{AutoGen (Microsoft)}
\label{subsec:autogen}

AutoGen pioneered the conversational multi-agent paradigm, where agents coordinate through natural language messages.

\textbf{Strengths:}
\begin{itemize}
    \item Flexible, emergent coordination
    \item Good for research and experimentation
    \item Strong code execution capabilities
    \item Active research community
\end{itemize}

\textbf{Limitations for Enterprise:}
\begin{itemize}
    \item Designed for research, not production
    \item No built-in persistence or recovery
    \item Emergent behavior difficult to predict/audit
    \item Minimal security/isolation features
\end{itemize}

\textbf{Recommendation:} AutoGen is suitable for research and prototyping but should not be used for production platform engineering without significant hardening.

\subsection{Framework Comparison Summary}
\label{subsec:framework-summary}

\textbf{Key insight:} No single framework addresses all enterprise requirements. Organizations should expect to build additional infrastructure around these tools.

\section{Integration with Platform Engineering Tooling}
\label{sec:platform-integration}

The value of agent-enhanced platform engineering depends on how well agents integrate with existing infrastructure tools.

\subsection{Kubernetes Operators}
\label{subsec:kubernetes-operators}

Kubernetes operators extend the Kubernetes API with custom resources and controllers. Agent integration can occur at multiple levels:

\textbf{Pattern A: Agent as Operator Consumer}

Agents interact with existing operators via standard Kubernetes APIs.

Implementation considerations:
\begin{itemize}
    \item Agent needs kubeconfig/service account with appropriate RBAC
    \item Read operations (get, list, watch) are safe for exploration
    \item Write operations (create, patch, delete) require careful scoping
    \item Agents can generate YAML and submit via GitOps
\end{itemize}

\textbf{Pattern B: Agent-Aware Operators}

Custom operators designed for agent interaction might expose:
\begin{itemize}
    \item Natural language intent in CRD specs
    \item Explanation fields in status
    \item Suggested actions for common issues
    \item Structured decision points for agent approval
\end{itemize}

\subsection{Terraform and OpenTofu}
\label{subsec:terraform}

Infrastructure as Code tools present multiple integration points:

\textbf{Pattern A: Agent-Generated IaC}

Agents generate Terraform/OpenTofu configurations from natural language.

Trust considerations:
\begin{itemize}
    \item Generated code MUST go through review (human or agent reviewer)
    \item Agents should not have terraform apply privileges
    \item Dry-run (terraform plan) can be agent-executed for validation
\end{itemize}

\textbf{Pattern B: Agent-Assisted Plan Analysis}

Agents analyze terraform plan output for risks:
\begin{itemize}
    \item Destructive operations (destroy, replace)
    \item Sensitive resource types (databases, secrets)
    \item Drift from expected state
    \item Cost implications
\end{itemize}

\subsection{GitOps (ArgoCD, Flux)}
\label{subsec:gitops}

GitOps provides a natural trust boundary for agent operations: \textbf{Git is the approval gate}.

\textbf{The GitOps Agent Pattern:}
\begin{itemize}
    \item \textbf{Agent:} Proposes changes (generates manifests, creates PRs)
    \item \textbf{Git:} Tracks all changes with attribution
    \item \textbf{Human:} Approves high-risk changes
    \item \textbf{GitOps controller:} Applies approved changes
\end{itemize}

This pattern enforces separation of concerns---agent NEVER touches production directly.

\subsection{CI/CD (GitHub Actions, GitLab CI)}
\label{subsec:cicd}

Agents can participate in CI/CD pipelines in several ways:

\textbf{Pattern A: Agent-Generated Workflows}

Agents create or modify workflow definitions.

\textbf{Pattern B: Agent as CI Job}

Agents run as steps within existing pipelines for code review, security analysis, or test generation.

\textbf{Pattern C: Agent-Triggered Pipelines}

Agents can trigger pipelines based on external events (e.g., automated rollback on critical alert).

\textbf{Trust Boundaries in CI/CD:}

\begin{table}[htbp]
\centering
\caption{CI/CD Trust Boundaries for Agents}
\label{tab:cicd-trust}
\begin{tabular}{lll}
\toprule
\textbf{Operation} & \textbf{Trust Level} & \textbf{Agent Capability} \\
\midrule
Generate workflow YAML & Medium & With review \\
Trigger workflow & High & With constraints \\
Access secrets & Critical & Never \\
Push to protected branches & Critical & Never \\
Approve deployments & Critical & Human only \\
\bottomrule
\end{tabular}
\end{table}

\section{Enterprise Deployment Patterns}
\label{sec:enterprise-patterns}

These patterns emerged from iterative stress-testing: proposing architectures, identifying failure modes, refining until the patterns survived contact with known enterprise constraints.

\subsection{Pattern: Agent-in-Toolbox}
\label{subsec:agent-toolbox}

\textbf{Context:} Large enterprises typically organize around teams with isolated tenants---each team has their own infrastructure, GitOps pipelines, and access management. Behind PAM, teams maintain ``toolboxes'' containing applications for manual infrastructure management.

\textbf{Problem:} Where do you deploy the agent?

\textbf{Solution:} Deploy the agent INTO the existing toolbox. The agent inherits the team's existing blast radius---it can affect exactly what the team can already affect manually, nothing more.

\textbf{Benefits:}
\begin{itemize}
    \item Zero new attack surface---agent uses existing access paths
    \item Blast radius pre-contained---same as team's existing permissions
    \item Operational simplicity---one thing to deploy, not new infrastructure per team
    \item Incremental adoption---teams can enable/disable agent features independently
\end{itemize}

\subsection{Pattern: Read-Only Default}
\label{subsec:readonly-default}

\textbf{Problem:} LLM-based agents are non-deterministic. How do you get value from agent assistance without accepting the risk of autonomous modifications?

\textbf{Solution:} Agents operate in read-only mode by default. They can observe, query, analyze, and recommend---but not modify. Write capabilities are opt-in per team, requiring explicit enablement.

Technical enforcement, not just prompting:
\begin{itemize}
    \item \textbf{Kubernetes:} RBAC with read-only ClusterRole bound to agent's ServiceAccount
    \item \textbf{Databases:} Dedicated user with SELECT-only grants
    \item \textbf{Cloud:} IAM policies with ReadOnly managed policies
    \item \textbf{Git:} Token with read permissions only
\end{itemize}

\textbf{The Trust Progression:}
\begin{lstlisting}[caption={Trust Progression Levels}]
Level 0: Read-only, all domains
Level 1: Read-only + write to non-production
Level 2: Read-only + write to production with approval
Level 3: Autonomous write within defined boundaries
Level 4: Full autonomy (rarely appropriate)
\end{lstlisting}

Most teams should stabilize at Level 1--2. Level 3+ requires demonstrated track record and robust guardrails.

\subsection{Pattern: Federated NOC Integration}
\label{subsec:federated-noc}

\textbf{Problem:} How does NOC get cross-tenant incident visibility without violating tenant isolation?

\textbf{Anti-Pattern:} NOC Super-Agent with read access to all tenants---privileged network paths to every tenant, single compromise = total visibility loss.

\textbf{Solution: Push-Based Federation}

Team agents push incident summaries OUT to NOC. NOC receives pre-digested analysis without direct tenant access.

\textbf{Benefits:}
\begin{itemize}
    \item Tenant isolation preserved---no privileged inbound access
    \item Audit clean---data flow is unidirectional, logged at push point
    \item Scalable---adding teams doesn't increase NOC access scope
    \item Degraded gracefully---if one team agent is down, others still report
\end{itemize}

\subsection{Pattern: Event-Driven Triage}
\label{subsec:event-driven}

\textbf{Problem:} When should an agent start investigating?

\textbf{Solution:} Monitoring webhooks trigger agent triage on threshold violations.

\begin{lstlisting}[caption={Event-Driven Triage Routing}]
Prometheus/Datadog/etc.
    |
    +- Info: Log only (no agent)
    |
    +- Warning: Webhook -> Agent (background triage)
    |                        -> Queue summary for team
    |
    +- Critical: Webhook -> Agent (immediate triage)
                             +- Summary to NOC
                             +- Enriched context to PagerDuty
\end{lstlisting}

\subsection{Pattern: Vault-Integrated Credential Management}
\label{subsec:vault-integration}

\textbf{Problem:} Agents need credentials to access systems. Hard-coded or long-lived credentials create security risk.

\textbf{Solution:} Integrate with HashiCorp Vault (or OpenBao) for dynamic, short-lived, scoped credentials.

\begin{lstlisting}[caption={Vault-Integrated Credential Flow}]
Agent Triage Session:
    |
    +- 1. Authenticate to Vault (AppRole / K8s auth)
    |
    +- 2. Request scoped credential
    |      -> "Give me read-only kubectl for namespace X"
    |
    +- 3. Receive short-lived credential (TTL: 15 min)
    |
    +- 4. Perform triage operations
    |
    +- 5. Session ends
    |
    +- 6. Credential auto-expires (or explicit revoke)
\end{lstlisting}

\textbf{Benefits:}
\begin{itemize}
    \item No long-lived secrets in agent configuration
    \item Credential scope enforced by Vault policy, not just agent prompting
    \item Audit trail of all credential access
    \item Compromise limited---attacker gets short-lived, scoped credential at most
\end{itemize}

\subsection{Pattern: Git-Backed Agent Memory}
\label{subsec:git-memory}

\textbf{Problem:} Agents with persistent memory accumulate team-specific knowledge. How do you persist, version, and govern agent memory?

\textbf{Solution:} Agent memory stored in Git repositories---version controlled, auditable, recoverable.

\begin{lstlisting}[caption={Git-Backed Memory Structure}]
team-alpha-agent-memory/
  +-- incidents/
  |     +-- 2026-01-15-db-pool-exhaustion.md
  |     +-- 2026-02-03-auth-cascade.md
  +-- learnings/
  |     +-- service-quirks.md
  |     +-- deployment-patterns.md
  +-- runbooks/
  |     +-- generated-from-incidents.md
  +-- MEMORY.md (active context)
  +-- .git/
\end{lstlisting}

\textbf{Benefits:}
\begin{itemize}
    \item Memory survives agent restarts, upgrades, migrations
    \item Version history enables ``what did agent believe when incident X happened?''
    \item Human review possible for significant learnings
    \item Backup is standard git workflow
\end{itemize}

\subsection{Pattern: ChatOps Delivery}
\label{subsec:chatops}

\textbf{Problem:} Teams live in Slack/Teams. Dashboards require context switching.

\textbf{Solution:} Agent posts triage summaries and incident updates to team channels.

\textbf{Why Output-Only:}

The ``@agent restart the pod'' trap:
\begin{itemize}
    \item If agent accepts commands in chat, users will issue them
    \item Chat context is ambiguous, easy to misinterpret
    \item No approval workflow in chat interface
    \item Audit trail unclear
\end{itemize}

Instead: Agent provides analysis in chat. If human wants agent to act, they go to toolbox UI where proper guardrails exist.

\subsection{Enterprise Pattern Summary}
\label{subsec:pattern-summary}

\begin{table}[htbp]
\centering
\caption{Enterprise Pattern Language Summary}
\label{tab:enterprise-patterns}
\begin{tabular}{ll}
\toprule
\textbf{Pattern} & \textbf{Core Principle} \\
\midrule
Agent-in-Toolbox & Inherit existing access boundaries, don't create new ones \\
Read-Only Default & Capability restriction as risk management \\
Federated NOC & Push model preserves isolation \\
Event-Driven Triage & Agents activate on events, not polling \\
Cascading Failure Coordination & Hierarchy prevents thundering herd \\
Vault Credentials & Dynamic, scoped, short-lived secrets \\
Git-Backed Memory & Version control for institutional knowledge \\
ChatOps Delivery & Output where teams work, commands where guardrails exist \\
\bottomrule
\end{tabular}
\end{table}

These patterns emerged from stress-testing proposals against real constraints. They are not theoretically optimal---they are practically survivable.

\section{Summary}
\label{sec:implementation-summary}

This chapter has examined the practical dimensions of implementing agent-enhanced platform engineering:

\begin{itemize}
    \item \textbf{Topology choices} that determine trust boundaries and operational characteristics
    \item \textbf{State management approaches} with different tradeoffs for reliability and auditability
    \item \textbf{Tool ecosystem} gaps that organizations must address
    \item \textbf{Integration patterns} with established platform engineering infrastructure
    \item \textbf{Enterprise deployment patterns} that survive contact with real organizational constraints
\end{itemize}

The trust framework from Chapter~\ref{ch:trust-framework} is scale-invariant---the principles apply whether you're a startup with a single jump host or an enterprise with 100+ teams. What changes is the \textit{implementation}, not the \textit{principle}.

\chapter{Evaluation and Limitations}
\label{ch:evaluation}

% Working draft - S. Vene, 2026-02-08

\section{Introduction}
\label{sec:evaluation-intro}

This thesis proposes a trust framework and operational patterns for agent-enhanced platform engineering. This chapter addresses how these contributions can be evaluated and honestly acknowledges their limitations.

\section{Validation Approach}
\label{sec:validation-approach}

\subsection{Design Science Evaluation}
\label{subsec:design-science-evaluation}

Following the Design Science Research paradigm \citep{hevner2004,peffers2007}, the contributions are evaluated against:

\textbf{Utility:} Do the artifacts solve the identified problems?
\begin{itemize}
    \item Trust framework: Does it help organizations reason about appropriate autonomy levels?
    \item Operational patterns: Do they provide actionable guidance for implementation?
\end{itemize}

\textbf{Quality:} Are the artifacts well-constructed?
\begin{itemize}
    \item Internal consistency: Do patterns and framework components align?
    \item Completeness: Are major use cases covered?
    \item Clarity: Can practitioners understand and apply the guidance?
\end{itemize}

\textbf{Efficacy:} Do the artifacts work in practice?
\begin{itemize}
    \item Practitioner review: Do experienced platform engineers find the patterns plausible?
    \item Constraint satisfaction: Do patterns survive stress-testing against real enterprise constraints?
\end{itemize}

\subsection{Validation Activities Conducted}
\label{subsec:validation-activities}

\begin{table}[htbp]
\centering
\caption{Validation Activities}
\label{tab:validation-activities}
\begin{tabular}{lll}
\toprule
\textbf{Activity} & \textbf{Purpose} & \textbf{Status} \\
\midrule
Literature grounding & Trust framework rooted in Mayer et al. (1995) & Complete \\
Pattern stress-testing & Enterprise patterns tested against constraints & Complete \\
Practitioner experience & Author's 15 years applied to development & Complete \\
Industry evidence collection & Documented deployments validating patterns & Ongoing \\
Peer review & Academic and practitioner feedback & Ongoing \\
Empirical validation & Deployment and measurement in real environments & Future work \\
\bottomrule
\end{tabular}
\end{table}

\subsection{Industry Evidence}
\label{subsec:industry-evidence}

While controlled empirical validation remains future work, emerging industry deployments provide practitioner-level validation that the patterns in this thesis reflect real-world practice.

\textbf{HolmesGPT / Robusta.dev (KubeCon 2026)}

Elroy Parkinson reported at KubeCon that production deployments of HolmesGPT achieve 80\% correct root cause diagnosis in under 2 minutes. Their architecture directly implements this thesis's trust framework:

\begin{table}[htbp]
\centering
\caption{HolmesGPT Trust Framework Implementation}
\label{tab:holmesgpt-trust}
\begin{tabular}{ll}
\toprule
\textbf{Trust Factor} & \textbf{HolmesGPT Implementation} \\
\midrule
Observability & ``Strong observability infrastructure'' as explicit prerequisite \\
Reversibility & PR approval gate---agent proposes, human approves \\
Blast Radius & Read-only access---no direct mutation of production \\
Autonomy & L2--L3: Agent performs diagnosis, human retains execution authority \\
\bottomrule
\end{tabular}
\end{table}

Their finding that ``AI agents struggle with poorly maintained environments'' validates this thesis's position that observability patterns (Chapter~\ref{ch:observability}) are prerequisites, not optional enhancements.

\textbf{Limitations of industry evidence:}
\begin{itemize}
    \item Vendor claims lack published methodology
    \item Selection bias toward success stories
    \item No controlled comparison with non-agent approaches
\end{itemize}

Industry evidence is cited as \textit{practitioner validation}---confirmation that the patterns reflect real deployments---not as \textit{empirical proof} of efficacy.

\subsection{What Has Been Validated}
\label{subsec:validated}

\textbf{Trust Framework (Chapter~\ref{ch:trust-framework}):}
\begin{itemize}
    \item Grounded in established organizational trust theory \citep{mayer1995}
    \item Components (Observability, Reversibility, Blast Radius, Autonomy) derived from operational practice
    \item Constraint-based approach is logically consistent
    \item Practitioner review confirms face validity
\end{itemize}

\textbf{Operational Patterns (Chapters~\ref{ch:observability}--\ref{ch:remediation}):}
\begin{itemize}
    \item Patterns derived from industry practice and literature
    \item Stress-tested against enterprise constraints (multi-tenancy, PAM, compliance)
    \item Aligned with emerging industry tools and approaches
\end{itemize}

\textbf{Enterprise Deployment Patterns (Chapter~\ref{ch:implementation}):}
\begin{itemize}
    \item Derived from design exploration against real constraints
    \item Address gaps identified in current tool ecosystem
    \item Cross-referenced with industry best practices (Zero Trust, PAM integration)
\end{itemize}

\section{Limitations}
\label{sec:limitations}

\subsection{Methodological Limitations}
\label{subsec:methodological-limitations}

\textbf{Practitioner Bias:}
The patterns emerge from the author's experience in government and financial services. Other sectors (healthcare, retail, startups) may face different constraints and opportunities. The patterns should be adapted, not adopted wholesale.

\textbf{Single Practitioner Perspective:}
While informed by 15 years of experience across multiple organizations, the thesis reflects one practitioner's viewpoint. Broader validation through multiple practitioners would strengthen generalizability.

\textbf{Design Science vs. Empirical Research:}
The thesis proposes artifacts (frameworks, patterns) rather than testing hypotheses about the world. This is appropriate for the current state of the field but means contributions are ``proposed solutions'' not ``proven solutions.''

\subsection{Validation Limitations}
\label{subsec:validation-limitations}

\textbf{No Controlled Experiments:}
The patterns have not been tested through randomized controlled trials comparing agent-enhanced vs. traditional approaches. Such experiments would require significant organizational commitment and time.

\textbf{No Longitudinal Deployment Data:}
Claims about patterns improving MTTR, reducing toil, or building trust are hypothesized, not measured. The original proposal anticipated longitudinal industrial case studies; this remains future work.

\textbf{Trust Framework Not Empirically Calibrated:}
The constraint thresholds (e.g., ``L4 requires automated rollback $<$15min'') represent practitioner judgment, not empirically derived cutoffs. Organizations should calibrate based on their risk tolerance.

\subsection{Scope Limitations}
\label{subsec:scope-limitations}

\textbf{Focus on Observability/Incident/Remediation:}
The thesis focuses on operational contexts (RQ1--3) rather than the full platform engineering lifecycle. Design-time and build-time agent applications are mentioned but not deeply explored.

\textbf{Enterprise Context:}
Patterns assume relatively mature platform engineering organizations. Startups or less mature organizations may need different approaches.

\textbf{Technology Snapshot:}
The LLM and agent ecosystem is evolving rapidly. Specific tool recommendations (LangGraph, CrewAI, etc.) may become outdated. The principles should remain applicable; implementation details may require updating.

\subsection{What This Thesis Does NOT Claim}
\label{subsec:not-claimed}

\begin{itemize}
    \item The Trust Framework is NOT a calculable formula producing definitive scores
    \item The patterns are NOT guaranteed to improve MTTR or reduce incidents
    \item The autonomy levels are NOT empirically validated cutoffs
    \item Agent-enhanced operations does NOT replace human judgment for novel or high-stakes situations
\end{itemize}

\section{Future Work}
\label{sec:future-work}

\subsection{Empirical Validation}
\label{subsec:empirical-validation}

\textbf{Longitudinal Deployment Studies:}
Deploy patterns in partner organizations and measure:
\begin{itemize}
    \item Mean Time to Resolution (MTTR) before/after
    \item Incident recurrence rates
    \item Engineer cognitive load (surveys)
    \item Trust calibration accuracy (appropriate reliance)
\end{itemize}

\textbf{Trust Framework Calibration:}
Survey practitioners on trust decisions to empirically ground:
\begin{itemize}
    \item Which components matter most?
    \item What thresholds are appropriate for different contexts?
    \item How does trust evolve over time with agent experience?
\end{itemize}

\subsection{Pattern Extension}
\label{subsec:pattern-extension}

\textbf{Design-Time Patterns:}
Extend coverage to infrastructure design, including:
\begin{itemize}
    \item Agent-assisted architecture review
    \item Cost optimization recommendations
    \item Security posture assessment
\end{itemize}

\textbf{Build-Time Patterns:}
Extend coverage to CI/CD integration:
\begin{itemize}
    \item IaC generation and review
    \item Test generation for infrastructure
    \item Dependency vulnerability assessment
\end{itemize}

\subsection{Tool Development}
\label{subsec:tool-development}

\textbf{Trust Dashboard:}
Tool that operationalizes the trust framework:
\begin{itemize}
    \item Assess current safeguards against autonomy level requirements
    \item Track agent performance metrics for trust calibration
    \item Recommend autonomy adjustments based on track record
\end{itemize}

\textbf{Pattern Implementation Library:}
Reference implementations of operational patterns:
\begin{itemize}
    \item Observability synthesis agents
    \item Triage automation workflows
    \item Runbook execution frameworks
\end{itemize}

\section{Contribution to Knowledge}
\label{sec:contribution}

Despite limitations, this thesis contributes:

\begin{enumerate}
    \item \textbf{First constraint-based trust framework for operational agents}---Grounded in organizational trust theory, actionable for practitioners

    \item \textbf{Systematic treatment of observability/triage/remediation patterns}---Consolidates emerging practice with structured guidance

    \item \textbf{Enterprise deployment patterns}---Addresses real constraints (PAM, multi-tenancy, cost, data sovereignty) that existing frameworks overlook

    \item \textbf{Honest acknowledgment of limitations}---Presents contributions as hypotheses for practice, not proven solutions
\end{enumerate}

The contributions are offered as a foundation for both practice and future research.

\section{Summary}
\label{sec:evaluation-summary}

This chapter has:
\begin{itemize}
    \item Described the validation approach appropriate for Design Science Research
    \item Honestly acknowledged methodological, validation, and scope limitations
    \item Outlined future work for empirical validation and extension
    \item Positioned the contributions as foundations for ongoing research and practice
\end{itemize}

\chapter{Conclusion}
\label{ch:conclusion}

% Working draft - 2026-02-07

\section{Introduction}
\label{sec:conclusion-intro}

Platform engineering stands at an inflection point. The transition from deterministic automation to agent-augmented systems represents the most significant shift in infrastructure management since Infrastructure as Code. This thesis has examined how AI agents can enhance platform engineering across the full lifecycle---from design through operations---while maintaining the trust and governance required for enterprise environments.

This concluding chapter synthesizes the findings, definitively answers the research questions, articulates implications for key audiences, and charts directions for future research.

\section{Summary of Contributions}
\label{sec:summary-contributions}

This thesis makes five primary contributions to the field of agent-enhanced platform engineering:

\subsection{Contribution 1: Lifecycle Mapping}
\label{subsec:lifecycle-mapping}

We provided the first systematic analysis of where AI agents add value across the complete platform engineering lifecycle. The lifecycle mapping reveals that agent opportunities are not uniformly distributed. The Operate and Build phases offer the most mature and immediate opportunities, while Deploy and Design phases require more sophisticated human-agent collaboration patterns.

\subsection{Contribution 2: Integration Patterns}
\label{subsec:integration-patterns}

We documented reusable patterns for human-agent collaboration in platform work:
\begin{itemize}
    \item \textbf{The Augmentation Spectrum}---A four-level model (Inform $\rightarrow$ Assist $\rightarrow$ Automate with Review $\rightarrow$ Autonomous) that moves beyond binary automation decisions
    \item \textbf{Task Classification Matrix}---Mapping cognitive complexity against risk level to determine appropriate collaboration modes
    \item \textbf{Anti-patterns and failure modes}---Common mistakes in agent integration and their remediation
\end{itemize}

\subsection{Contribution 3: The Trust Framework}
\label{subsec:trust-framework-contribution}

We proposed a comprehensive framework for reasoning about trust in agentic systems, centered on the Trust Equation:

\begin{equation}
\text{Effective Trust} = \frac{\text{Observability} \times \text{Reversibility} \times \text{Blast Radius Control}}{\text{Autonomy Level}}
\end{equation}

Supporting constructs include:
\begin{itemize}
    \item \textbf{The Deterministic-Probabilistic Shift}---Articulating why traditional automation assumptions fail for agentic systems
    \item \textbf{Progressive Autonomy Model}---Adapting the ATF's ``earned trust'' concept for platform engineering contexts
    \item \textbf{Trust-Building Patterns}---Five concrete patterns (Audit Everything, Rewind Capability, Blast Radius Containment, Human Checkpoints, Progressive Trust)
    \item \textbf{The Complexity Cliff}---Analysis of how trust requirements scale non-linearly from single-team to enterprise deployments
\end{itemize}

\subsection{Contribution 4: Practical Guidance}
\label{subsec:practical-guidance}

We grounded all contributions in real-world practice, drawing on:
\begin{itemize}
    \item Implementation experience with multi-agent orchestration systems
    \item Case studies from production deployments
    \item Integration patterns with existing platform tooling (Kubernetes, Terraform, GitOps)
    \item Architectural options comparing session-based versus in-memory orchestration
\end{itemize}

\subsection{Contribution 5: Adoption Roadmap}
\label{subsec:adoption-roadmap}

We developed a structured approach for organizations to introduce agent-enhanced platform engineering:
\begin{itemize}
    \item \textbf{Fit Assessment Framework}---Six-criteria evaluation with scoring methodology
    \item \textbf{Progressive Advancement Criteria}---Measurable thresholds for increasing agent autonomy
    \item \textbf{Organizational Change Considerations}---Addressing skill gaps, team evolution, and building the business case
\end{itemize}

\section{Answering the Research Questions}
\label{sec:answering-rqs}

\subsection{RQ1: How can agentic systems enhance observability?}
\label{subsec:rq1-answer}

\textbf{Answer:} Agents add the most immediate value in the \textbf{Operate} and \textbf{Build} phases.

Specifically:
\begin{itemize}
    \item \textbf{Operate:} Alert triage and noise reduction can recover 2--4 hours per engineer per day. Root cause analysis assistance can reduce MTTR by 30--50\%.
    \item \textbf{Build:} IaC generation, security scanning, and test generation benefit from agent consistency and breadth of knowledge.
\end{itemize}

Effective integration follows the \textbf{Augmentation Spectrum}: start at Inform level, progress to Assist, then Automate with Review, and finally Autonomous---only for proven, low-risk tasks.

\subsection{RQ2: To what extent can agentic reasoning improve incident triage?}
\label{subsec:rq2-answer}

\textbf{Answer:} Enterprise-scale deployment requires four interlocking mechanisms:

\begin{enumerate}
    \item \textbf{Observability Infrastructure}---Comprehensive logging with decision provenance
    \item \textbf{Reversibility by Design}---Agent actions must be reversible by default
    \item \textbf{Blast Radius Containment}---Strict scoping through identity, network isolation, and resource limits
    \item \textbf{Progressive Autonomy Systems}---Treating agents like employees who earn trust through track record
\end{enumerate}

At enterprise scale, \textbf{federation} is essential: central framework with distributed implementation.

\subsection{RQ3: Under what conditions can agents perform remediation?}
\label{subsec:rq3-answer}

\textbf{Answer:} Collaboration structure should match task characteristics according to two dimensions:

\textbf{By Cognitive Complexity:}
\begin{itemize}
    \item \textbf{Procedural tasks:} Automate fully
    \item \textbf{Analytical tasks:} Agent-assisted with human review
    \item \textbf{Diagnostic tasks:} Human-led with agent support
    \item \textbf{Strategic tasks:} Human decision with agent information
\end{itemize}

\textbf{By Risk Level:}
\begin{itemize}
    \item \textbf{Low risk:} Higher autonomy acceptable
    \item \textbf{Medium risk:} Automate with review
    \item \textbf{High risk:} Human approval required
    \item \textbf{Critical risk:} Human-only with agent assistance
\end{itemize}

Key principle: \textbf{Agents and humans have complementary strengths.} Effective collaboration leverages both.

\subsection{RQ4: What trust and governance mechanisms enable appropriate autonomy?}
\label{subsec:rq4-answer}

\textbf{Answer:} Four primary barriers emerge, each with specific remediation:

\begin{enumerate}
    \item \textbf{Trust Deficit}---Address through progressive autonomy, starting with low-risk applications
    \item \textbf{Governance Gaps}---Select tools with enterprise focus, build governance layers
    \item \textbf{Skill Gaps}---Treat initial projects as learning investments
    \item \textbf{Regulatory Uncertainty}---Conservative interpretation, comprehensive audit trails
\end{enumerate}

The most effective adoption approach is \textbf{start small, prove value, expand systematically}.

\section{Implications and Takeaways}
\label{sec:implications}

\subsection{For Platform Engineers}
\label{subsec:for-engineers}

\begin{enumerate}
    \item \textbf{Start with operations, not deployment.} Alert triage and log analysis offer immediate value with bounded risk.
    \item \textbf{Think augmentation, not replacement.} Strategic decisions, stakeholder relationships, and novel problem-solving remain fundamentally human.
    \item \textbf{Build observability first.} Before deploying any agent, ensure you can answer: What did it do? Why? Can we undo it?
    \item \textbf{Embrace progressive autonomy.} Don't grant agents full access on day one.
    \item \textbf{Document everything.} Agentic systems require more rigorous documentation than traditional automation.
\end{enumerate}

\subsection{For Tool Developers}
\label{subsec:for-developers}

\begin{enumerate}
    \item \textbf{Enterprise-grade governance is not optional.}
    \item \textbf{Observability is a first-class requirement.}
    \item \textbf{Design for reversibility.}
    \item \textbf{Session-based orchestration deserves attention.}
    \item \textbf{Support progressive autonomy.}
\end{enumerate}

\subsection{For Organizations}
\label{subsec:for-organizations}

\begin{enumerate}
    \item \textbf{Agent adoption is a change management challenge, not just a technical one.}
    \item \textbf{Start with a trust framework.}
    \item \textbf{Invest in the adoption roadmap.}
    \item \textbf{Regulatory engagement is strategic.}
    \item \textbf{Measure what matters.}
\end{enumerate}

\subsection{For Researchers}
\label{subsec:for-researchers}

\begin{enumerate}
    \item \textbf{The field needs empirical validation.}
    \item \textbf{Session-based orchestration is underexplored.}
    \item \textbf{Human-agent collaboration metrics are underdeveloped.}
    \item \textbf{Longitudinal studies are essential.}
    \item \textbf{Regulatory implications require interdisciplinary work.}
\end{enumerate}

\section{Future Research Directions}
\label{sec:future-directions}

\subsection{Direction 1: Empirical Validation of the Trust Equation}
\label{subsec:direction1}

The Trust Equation provides a conceptual model, but empirical validation is needed to establish:
\begin{itemize}
    \item Whether the multiplicative relationship holds in practice
    \item Relative weighting of components across contexts
    \item Threshold values for effective deployment
    \item Organizational and cultural factors that moderate the relationship
\end{itemize}

\subsection{Direction 2: Standardized Human-Agent Collaboration Metrics}
\label{subsec:direction2}

The field lacks standardized metrics for measuring human-agent collaboration quality. Research is needed to develop and validate:
\begin{itemize}
    \item \textbf{Collaboration Efficiency Index}---Are agents helping humans work faster?
    \item \textbf{Decision Quality Score}---Are human-agent decisions better than either alone?
    \item \textbf{Escalation Appropriateness Rate}---Do agents escalate at the right times?
    \item \textbf{Trust Calibration Measure}---Is human trust aligned with agent capability?
\end{itemize}

\subsection{Direction 3: Context-Adaptive Autonomy Systems}
\label{subsec:direction3}

Current progressive autonomy models use static promotion criteria. Future systems should adapt autonomy levels dynamically based on:
\begin{itemize}
    \item Current context (incident in progress, maintenance window, etc.)
    \item Task characteristics (novel situation, routine operation)
    \item Agent confidence levels
    \item Environmental risk factors
\end{itemize}

\subsection{Direction 4: Multi-Agent Coordination Patterns}
\label{subsec:direction4}

As deployments scale, multiple agents must coordinate. Research is needed on:
\begin{itemize}
    \item Coordination protocols that maintain trust guarantees
    \item Conflict resolution when agents disagree
    \item Collective responsibility and accountability attribution
    \item Emergent behavior detection and prevention in multi-agent systems
\end{itemize}

\subsection{Direction 5: Regulatory Framework Development}
\label{subsec:direction5}

The intersection of AI governance frameworks (AI Act, NIS2) and agentic systems requires research on:
\begin{itemize}
    \item How existing regulations apply to autonomous platform agents
    \item What additional requirements may be needed for agentic systems
    \item How organizations can demonstrate compliance for probabilistic systems
    \item Standards for agent certification and audit
\end{itemize}

\section{Limitations}
\label{sec:conclusion-limitations}

This thesis has several limitations that should inform interpretation:

\textbf{Scope:} The focus on platform engineering, while enabling depth, limits generalizability to other domains.

\textbf{Empirical Grounding:} While grounded in real implementation experience and industry practice, the proposed frameworks lack rigorous empirical validation through controlled studies.

\textbf{Temporal Sensitivity:} The field is evolving rapidly. LLM capabilities, tooling ecosystems, and regulatory frameworks will change significantly during the life of this work.

\textbf{Organizational Context:} Recommendations assume relatively mature platform engineering organizations.

\textbf{Tool Ecosystem Focus:} The analysis references specific tools that may be superseded. The principles should remain applicable, but specific guidance may require updating.

\section{Final Reflections: The Road Ahead}
\label{sec:final-reflections}

Platform engineering is evolving from \textbf{infrastructure automation} to \textbf{infrastructure augmentation}. The shift from deterministic pipelines to probabilistic agents represents not just a technical change but a fundamental reconceptualization of how humans and machines collaborate on infrastructure.

This transition will not be without friction. Engineers trained in deterministic systems must learn to work with probabilistic ones. Organizations accustomed to rule-based governance must develop new mechanisms for autonomous agents. Regulators struggling to understand AI in general must grapple with domain-specific applications.

Yet the trajectory is clear. The pain points in platform engineering---complexity, documentation debt, alert fatigue, expertise bottlenecks---are not solvable through more deterministic automation. The marginal returns on traditional approaches are diminishing. Agentic augmentation offers a path to the next level of platform capability.

The key insight from this research is that \textbf{agent integration is not binary}. The choice is not ``automate or don't'' but rather ``how much, for what tasks, with what safeguards.'' The frameworks presented here---the Trust Equation, the Augmentation Spectrum, the Fit Assessment---provide tools for making these nuanced decisions.

Success will come to organizations that:
\begin{enumerate}
    \item \textbf{Start conservatively}---building trust through demonstration, not declaration
    \item \textbf{Invest in governance}---recognizing that trust infrastructure is not overhead but enabler
    \item \textbf{Treat agents as team members}---with appropriate onboarding, supervision, and progression
    \item \textbf{Measure what matters}---developing new metrics for human-agent collaboration
    \item \textbf{Learn continuously}---adapting as capabilities and requirements evolve
\end{enumerate}

The future of platform engineering is not AI replacing humans or humans rejecting AI. It is humans and agents working together, each contributing their strengths, governed by frameworks that enable trust while preserving control.

That future is being built now. This thesis contributes to its foundations.


% Back matter
\backmatter

% Acknowledgements
\chapter*{Acknowledgements}
\addcontentsline{toc}{chapter}{Acknowledgements}

[To be written]

% Bibliography
\bibliographystyle{apalike}
\bibliography{references}

% Abstracts
\begin{abstracten}
This thesis develops a trust framework and operational patterns for integrating AI agents into platform engineering, with focus on observability, incident triage, and remediation. The core contribution is a constraint-based trust model grounded in organizational trust theory (Mayer et al., 1995) that enables organizations to calibrate appropriate autonomy levels for operational agents.

The research addresses four questions: (1) How can agents enhance observability? (2) Can agentic reasoning improve incident triage? (3) Under what conditions can agents safely perform remediation? (4) What trust mechanisms enable appropriate autonomy calibration?

Using Design Science methodology grounded in 15 years of platform engineering practice, this thesis proposes: a Trust Equation mapping Observability, Reversibility, and Blast Radius to autonomy levels; operational patterns for agent-enhanced observability, triage, and remediation; and governance guidance for enterprise deployment.

The work contributes to the emerging field of agent-enhanced operations by providing actionable frameworks for practitioners seeking to deploy AI agents in production systems responsibly.

\textbf{Keywords:} AI agents, platform engineering, trust framework, observability, incident response, autonomous systems, LLM
\end{abstracten}

\begin{abstractet}
Käesolev doktoritöö arendab usaldusraamistikku ja operatiivseid mustreid tehisintellekti agentide integreerimiseks platvormiarendusse, keskendudes jälgitavusele, intsidentide triaazhile ja remediatsioonile. Peamine panus on piirangupõhine usaldus\-mudel, mis põhineb organisatsioonilise usalduse teoorial (Mayer et al., 1995) ja võimaldab organisatsioonidel kalibreerida operatiivsete agentide jaoks sobivaid autonoomia tasemeid.

Uurimistöö käsitleb nelja küsimust: (1) Kuidas saavad agendid parandada jälgitavust? (2) Kas agentne arutlus parandab intsidentide triaazhit? (3) Millistel tingimustel saavad agendid ohutult teostada remediatsiooni? (4) Millised usaldusmehhanismid võimaldavad sobivat autonoomia kalibreerimist?

Kasutades disainiteaduse metoodikat, mis põhineb 15-aastasel platvormiarenduse praktikal, pakub käesolev töö: usaldusvõrrandit, mis seob jälgitavuse, pööratavuse ja mõjuulatuse autonoomia tasemetega; operatiivseid mustreid agentidega täiustatud jälgitavuse, triazhi ja remediatsiooni jaoks; ning juhtimissuuniseid ettevõtte juurutamiseks.

\textbf{Võtmesõnad:} tehisintellekti agendid, platvormiarendus, usaldusraamistik, jälgitavus, intsidentidele reageerimine, autonoomsed süsteemid, suur keele\-mudel
\end{abstractet}

% Curriculum Vitae
\begin{curriculumvitae}
\textbf{Siim Vene}

\section*{Personal Information}
Date of birth: January 21, 1980\\
Citizenship: Estonian

\section*{Contact}
Email: siim.vene@taltech.ee

\section*{Education}
\begin{tabular}{ll}
2026-- & PhD (ongoing), TalTech, Information Technology \\
2017 & MSc, TalTech, ``Implementing DevOps in large organizations'' \\
2001--2005 & Associate Degree, Estonian IT College, IT Systems Administration \\
\end{tabular}

\section*{Professional Experience}
\begin{tabular}{ll}
2020-- & Chief Infrastructure Architect, SMIT \\
2019-- & Program Manager, TalTech \\
2023-- & Public Cloud Service Manager, Riigi IT Keskus \\
2025-- & Founder, Kleidia \\
2017--2020 & Domain Architect, Swedbank \\
\end{tabular}

\section*{Research Interests}
Platform engineering, AI agents, observability, incident response, trust frameworks, autonomous systems

\section*{Publications}
[To be updated]
\end{curriculumvitae}

\begin{elulookirjeldus}
\textbf{Siim Vene}

\section*{Isikuandmed}
Sünniaeg: 21. jaanuar 1980\\
Kodakondsus: Eesti

\section*{Kontakt}
E-post: siim.vene@taltech.ee

\section*{Haridus}
\begin{tabular}{ll}
2026-- & Doktorantuur (pooleli), TalTech, Infotehnoloogia \\
2017 & MSc, TalTech, ``Implementing DevOps in large organizations'' \\
2001--2005 & Rakenduskõrgharidus, Eesti Infotehnoloogia Kolledž, IT-süsteemide haldus \\
\end{tabular}

\section*{Töökogemus}
\begin{tabular}{ll}
2020-- & Taristu peaarhitekt, SMIT \\
2019-- & Programmijuht, TalTech \\
2023-- & Avaliku pilve teenusehaldur, Riigi IT Keskus \\
2025-- & Asutaja, Kleidia \\
2017--2020 & Domeeniarhitekt, Swedbank \\
\end{tabular}

\section*{Teadustöö huvialad}
Platvormiarendus, tehisintellekti agendid, jälgitavus, intsidentidele reageerimine, usaldusraamistikud, autonoomsed süsteemid
\end{elulookirjeldus}

\end{document}
